\chapter{Coequational Subrecognizers and Syntactic Transition--Output Images}
\label{ch:mealy-coequations-syntactic-images}

\section{Overview}
\label{sec:coeq-overview}

The preceding chapter constructed, for a behaviour specification
\[
    k:K\to \mathbf{Beh}_0,
\]
its derivative image
\[
    \operatorname{Der}_0(K)\hookrightarrow \mathbf{Beh}_0
\]
and the corresponding minimal residual Mealy morphism
\[
    \mathsf{Min}(K):\mathbb A\rightsquigarrow\mathbb B.
\]
The purpose of the present chapter is to reorganize this construction in
coequational form and then to associate to every such coequational theory a
syntactic transition--output category.

The central semantic object is the canonical behaviour recognizer
\[
    \mathbf{Beh}^{\mathrm{can}}_{\mathbb A,\mathbb B}:
    \mathbb A\rightsquigarrow\mathbb B.
\]
By Chapter~\ref{ch:behaviour-categories-myhill-nerode}, it is terminal among
strict residual semantic recognizers. Thus a local coequational theory is,
categorically, a subobject of this terminal recognizer:
\[
    V\hookrightarrow \mathbf{Beh}^{\mathrm{can}}_{\mathbb A,\mathbb B}.
\]
A recognizer
\[
    P:\mathbb A\rightsquigarrow\mathbb B
\]
models \(V\) precisely when its unique semantic map to the terminal recognizer
factors through \(V\):
\[
\begin{tikzcd}[column sep=large, row sep=large]
    P
        \arrow[rr, "\beta_P"]
        \arrow[dr, dashed, "\bar\beta_P"']
    &&
    \mathbf{Beh}^{\mathrm{can}}_{\mathbb A,\mathbb B}
    \\
    &
    V
        \arrow[ur, hook, "j_V"']
    &
\end{tikzcd}
\]
Equivalently, on carriers, the residual semantics
\[
    \beta_P:P\to\mathbf{Beh}_0
\]
factors through the underlying subobject
\[
    V\hookrightarrow\mathbf{Beh}_0.
\]

The basic examples of coequational theories are semantic images. For a
recognizer \(P\), its semantic image is the regular image of its terminal
semantic map:
\[
    P\twoheadrightarrow \operatorname{Beh}(P)
    \hookrightarrow \mathbf{Beh}^{\mathrm{can}}_{\mathbb A,\mathbb B}.
\]
Since regular images of strict recognizer maps are created by the ambient slice
category, \(\operatorname{Beh}(P)\) is again a subrecognizer of the canonical
behaviour recognizer.

The chapter proves a coequational Birkhoff theorem. Locally, a replete full
class of recognizers is definable by a local coequational theory if and only if
it is closed under strict homomorphic preimages, regular quotient recognizers,
and small coproducts. The representing coequational theory associated to a
class
\[
    \mathcal C\subseteq\mathsf{Rec}_{\mathbb A,\mathbb B}
\]
is the join of the semantic images of its recognizers:
\[
    V_{\mathcal C}
    :=
    \bigvee_{P\in\mathcal C}\operatorname{Beh}(P)
    \hookrightarrow \mathbf{Beh}_0.
\]
The use of joins of subobjects is the point at which this chapter strengthens
the ambient assumption to a Grothendieck topos.

If the output category carries algebraic structure, the same subobject
semantics gives the structured version. When \(\mathbb B\) is an algebra
object in internal categories for a finite-product theory \(\mathbb T\), the
behaviour object \(\mathbf{Beh}_0\) inherits pointwise \(\mathbb T\)-operations.
A local \(\mathbb T\)-coequational theory is a local coequational theory closed
under these pointwise operations. On model classes this corresponds to closure
under the induced \(\mathbb T\)-products of recognizers.

Globally, when the input category varies, an internal functor
\[
    F:\mathbb A'\to\mathbb A
\]
induces input pullback of recognizers
\[
    F^*:\mathsf{Rec}_{\mathbb A,\mathbb B}
    \to
    \mathsf{Rec}_{\mathbb A',\mathbb B}
\]
and residual reindexing of behaviours
\[
    F^\sharp:
    A'_0\times_{F_0,A_0,p_{\mathbf{Beh}}}
    \mathbf{Beh}_{\mathbb A,\mathbb B,0}
    \to
    \mathbf{Beh}_{\mathbb A',\mathbb B,0}.
\]
The key semantic equation is
\[
    \beta_{F^*P}=F^\sharp\beta_P.
\]
Thus global coequational theories are precisely local theories stable under
these behaviour reindexing maps.

The main syntactic construction of the chapter is the operation image of a
coequational theory. A Mealy recognizer with carrier
\[
    A_0\xleftarrow{p}P\xrightarrow{q}B_0
\]
can be equivalently viewed as a fibrewise Kleisli action
\[
    \chi_P:
    \mathbb A^{\mathrm{op}}
    \longrightarrow
    \operatorname{KlEnd}_{\mathbb B}(P),
\]
where \(\operatorname{KlEnd}_{\mathbb B}(P)\) is the internal category of
fibrewise output-comparison Kleisli endomorphisms of \(P\). Under this
correspondence, an arrow \(a:u\to v\) of \(\mathbb A\), regarded as
\(a^{\mathrm{op}}:v\to u\), acts by
\[
    x\longmapsto
    \bigl(\delta_P(a,x),\omega_P(a,x)\bigr).
\]
For a local coequational theory
\[
    V\hookrightarrow \mathbf{Beh}^{\mathrm{can}},
\]
we apply this to the restricted canonical recognizer on \(V\). Its action is
\[
    \chi_V:
    \mathbb A^{\mathrm{op}}
    \longrightarrow
    \operatorname{KlEnd}_{\mathbb B}(V).
\]
The transition--output image, or operation image, of \(V\) is the
identity-on-objects regular image of this action:
\[
    \operatorname{Op}_{\mathbb A,\mathbb B}(V)
    :=
    \operatorname{RegIm}(\chi_V).
\]
Thus
\[
    \operatorname{Op}_{\mathbb A,\mathbb B}(V)
\]
is the quotient of \(\mathbb A^{\mathrm{op}}\) by equality of the
transition--output operations induced on the states of \(V\).

For a specification
\[
    k:K\to\mathbf{Beh}_0,
\]
the syntactic transition--output category is
\[
    \operatorname{SynTO}_{\mathbb A,\mathbb B}(K)
    :=
    \operatorname{Op}_{\mathbb A,\mathbb B}(\operatorname{Der}_0(K)).
\]
It is the operation image of the smallest residual coequational theory
generated by the specification.

This viewpoint is close to coequational formulations of Birkhoff- and
Eilenberg-type theorems, where language-side theories are stable subobjects of
a canonical behaviour object and algebraic recognizers are systems whose
semantics land in those subobjects
\cite{Salamanca2017Unveiling,Salamanca2016Monad}. The finite and predual
perspective on generalized Eilenberg correspondences is also in the background
\cite{AdamekMiliusMyersUrbat2015Local,AdamekMiliusMyersUrbat2019TOCL}.

\begin{assumption}
\label{ass:coeq-grothendieck-topos}
    Throughout this chapter, \(\mathcal E\) is a Grothendieck topos with chosen
    pullbacks. Thus \(\mathcal E\) is Barr-exact and locally cartesian closed,
    has stable regular epi--mono factorizations, effective quotients of internal
    equivalence relations, small coproducts, and complete subobject lattices in
    every slice. All coproducts and joins below are small, relative to the fixed
    ambient universe.
\end{assumption}

The input and output internal categories are
\[
    \mathbb A=(A_0,A_1,s_A,t_A,e_A,m_A),
    \qquad
    \mathbb B=(B_0,B_1,s_B,t_B,e_B,m_B).
\]
When the input category is fixed, we write
\[
    Z:=A_0\times B_0.
\]
Composition in \(\mathbb A\) is written
\[
    m_A(a,b)=b\circ a
\]
for composable arrows
\[
    a:u\to v,
    \qquad
    b:v\to w.
\]
Thus a Mealy state over \(w\) processes \(b\circ a\) by first processing \(b\),
then processing \(a\).

\section{Coequational recognizers}
\label{sec:coequational-recognizers}

This section isolates the strict recognizer category, identifies local
coequational theories as subobjects of its terminal semantic object, and proves
the local coequational Birkhoff theorem in this setting.

\subsection{Recognizers and strict homomorphisms}
\label{subsec:recognizers-strict-homomorphisms}

\begin{definition}[Mealy recognizer]
\label{def:mealy-recognizer}
    A \textbf{Mealy recognizer} over \((\mathbb A,\mathbb B)\) is a Mealy
    morphism
    \[
        P:\mathbb A\rightsquigarrow\mathbb B
    \]
    with carrier span
    \[
        A_0\xleftarrow{p_P}P\xrightarrow{q_P}B_0.
    \]
    Thus \(P\) consists of transition and output maps
    \[
        \delta_P:
        A_1\times_{t_A,A_0,p_P}P\to P
    \]
    and
    \[
        \omega_P:
        A_1\times_{t_A,A_0,p_P}P\to B_1
    \]
    satisfying the Mealy unit and multiplication laws. In this chapter the word
    ``recognizer'' means a Mealy morphism considered through its residual
    semantics.
\end{definition}

\begin{definition}[Strict homomorphism]
\label{def:strict-homomorphism}
    Let
    \[
        P,Q:\mathbb A\rightsquigarrow\mathbb B
    \]
    be Mealy recognizers. A \textbf{strict homomorphism}
    \[
        f:P\to Q
    \]
    is a map in \(\mathcal E/Z\), so
    \[
        p_Qf=p_P,
        \qquad
        q_Qf=q_P,
    \]
    satisfying
    \[
        f(\delta_P(a,x))=\delta_Q(a,fx)
    \]
    and
    \[
        \omega_P(a,x)=\omega_Q(a,fx)
    \]
    for every generalized element
    \[
        (a,x)\in A_1\times_{t_A,A_0,p_P}P.
    \]
    We write
    \[
        \mathsf{Rec}_{\mathbb A,\mathbb B}
    \]
    for the category of Mealy recognizers and strict homomorphisms.
\end{definition}

The strict homomorphism equations are summarized by
\[
\begin{tikzcd}[column sep=large, row sep=large]
    x
        \arrow[r, "{a,\omega_P(a,x)}"]
        \arrow[d, "f"']
    &
    \delta_P(a,x)
        \arrow[d, "f"]
    \\
    fx
        \arrow[r, "{a,\omega_Q(a,fx)}"']
    &
    \delta_Q(a,fx).
\end{tikzcd}
\]

\begin{definition}[Subrecognizer and regular quotient recognizer]
\label{def:sub-and-quotient-recognizer}
    A \textbf{subrecognizer} is a monic strict homomorphism
    \[
        Q\hookrightarrow P.
    \]
    A \textbf{regular quotient recognizer} is a regular epimorphic strict
    homomorphism
    \[
        P\twoheadrightarrow Q.
    \]
    A class of recognizers is \textbf{closed under strict homomorphic
    preimages} if, for every strict homomorphism
    \[
        f:P\to Q,
    \]
    membership of \(Q\) implies membership of \(P\).
\end{definition}

\begin{remark}[Strict homomorphisms versus simulations]
\label{rem:strict-homomorphisms-versus-simulations}
    Strict homomorphisms should not be confused with the Mealy simulations of
    Chapter~\ref{ch:spans-internal-categories-mealy}. A simulation
    \[
        P\Rightarrow Q
    \]
    has an oppositely directed state map
    \[
        Q\to P
    \]
    and compares outputs by arrows of \(\mathbb B\). A strict homomorphism
    \[
        P\to Q
    \]
    is instead an ordinary map in \(\mathcal E/Z\) preserving transitions and
    outputs on the nose.
\end{remark}

Recall from Chapter~\ref{ch:behaviour-categories-myhill-nerode} that every
recognizer \(P\) has a residual semantics map
\[
    \beta_P:P\to\mathbf{Beh}_0.
\]
For \(x\in P\), the residual behaviour
\[
    \beta_P(x)=H_x:\mathbb A/p_P(x)\to\mathbb B
\]
is given on objects by
\[
    H_x(r)=q_P(\delta_P(r,x)),
\]
and on residual arrows
\[
    w\xrightarrow{a}u\xrightarrow{r}p_P(x)
\]
by
\[
    H_x(a,r)=\omega_P(a,\delta_P(r,x)).
\]

\begin{lemma}[Strict homomorphisms preserve residual semantics]
\label{lem:strict-homs-preserve-beta}
    Let
    \[
        f:P\to Q
    \]
    be a strict homomorphism. Then
    \[
        \beta_Qf=\beta_P.
    \]
\end{lemma}

\begin{proof}
    \leavevmode
    \begin{enumerate}
        \item \textbf{Compare object assignments.}
        Let \(x\in P\), and let
        \[
            r:u\to p_P(x)
        \]
        be a residual object. Since \(f\) preserves transitions and lies over
        \(B_0\),
        \[
        \begin{aligned}
            q_P(\delta_P(r,x))
            &=
            q_Q(f(\delta_P(r,x))) \\
            &=
            q_Q(\delta_Q(r,fx)).
        \end{aligned}
        \]
        Thus \(\beta_P(x)\) and \(\beta_Q(fx)\) agree on residual objects.

        \item \textbf{Compare arrow assignments.}
        For a residual arrow
        \[
            w\xrightarrow{a}u\xrightarrow{r}p_P(x),
        \]
        strictness gives
        \[
            \delta_Q(r,fx)=f(\delta_P(r,x)),
        \]
        and hence
        \[
        \begin{aligned}
            \omega_Q(a,\delta_Q(r,fx))
            &=
            \omega_Q(a,f(\delta_P(r,x))) \\
            &=
            \omega_P(a,\delta_P(r,x)).
        \end{aligned}
        \]
        Thus the two residual behaviours agree on arrows.

        \item \textbf{Display the semantic square.}
        The equality is the commutativity of
        \[
        \begin{tikzcd}[column sep=large, row sep=large]
            P
                \arrow[r, "f"]
                \arrow[dr, "\beta_P"']
            &
            Q
                \arrow[d, "\beta_Q"]
            \\
            &
            \mathbf{Beh}_0 .
        \end{tikzcd}
        \]
    \end{enumerate}
\end{proof}

\begin{theorem}[Terminal residual semantic recognizer]
\label{thm:terminal-recognizer-coeq}
    The canonical behaviour recognizer
    \[
        \mathbf{Beh}^{\mathrm{can}}_{\mathbb A,\mathbb B}
    \]
    is terminal in
    \[
        \mathsf{Rec}_{\mathbb A,\mathbb B}.
    \]
    For every recognizer \(P\), the unique strict homomorphism
    \[
        P\to\mathbf{Beh}^{\mathrm{can}}_{\mathbb A,\mathbb B}
    \]
    is its residual semantics map
    \[
        \beta_P.
    \]
\end{theorem}

\begin{proof}
    \leavevmode
    \begin{enumerate}
        \item \textbf{Use terminal residual semantics.}
        This is exactly Theorem~\ref{thm:terminal-residual-semantics}, expressed
        in the recognizer terminology of the present chapter.

        \item \textbf{Use uniqueness of strict maps to the canonical recognizer.}
        If \(f:P\to\mathbf{Beh}^{\mathrm{can}}\) is strict, then the uniqueness
        argument of Chapter~\ref{ch:behaviour-categories-myhill-nerode}
        reconstructs the value \(f(x)\) as the residual behaviour \(H_x\). Hence
        \(f=\beta_P\).
    \end{enumerate}
\end{proof}

\subsection{Local coequational theories}
\label{subsec:local-coequational-theories}

The canonical behaviour object
\[
    \mathbf{Beh}_0
\]
carries the canonical behaviour recognizer
\[
    \mathbf{Beh}^{\mathrm{can}}_{\mathbb A,\mathbb B}:
    \mathbb A\rightsquigarrow\mathbb B.
\]
A state over \(v\in A_0\) is a residual behaviour
\[
    H:\mathbb A/v\to\mathbb B.
\]
For an arrow
\[
    a:u\to v,
\]
the canonical transition is the derivative
\[
    \partial_aH=H\circ a_*,
\]
and the canonical output comparison is
\[
    \omega_{\mathbf{Beh}}(a,H):
    q_{\mathbf{Beh}}(\partial_aH)\to q_{\mathbf{Beh}}(H),
\]
the arrow assigned by \(H\) to the residual morphism
\[
    a\to 1_v
\]
in \(\mathbb A/v\).

The output comparison is displayed by applying \(H\) to the residual triangle
\[
\begin{tikzcd}[column sep=large, row sep=large]
    u
        \arrow[r, "a"]
        \arrow[rr, "a", bend left=15]
    &
    v
        \arrow[r, "1_v"']
    &
    v .
\end{tikzcd}
\qquad
\leadsto
\qquad
\begin{tikzcd}[column sep=large, row sep=large]
    H(a)
        \arrow[r, "{H(a,1_v)}"]
    &
    H(1_v).
\end{tikzcd}
\]

Throughout this chapter, derivative domains are regarded over
\[
    Z=A_0\times B_0
\]
by the derivative state. Thus if
\[
    a:u\to v
\]
and \(H\) lies over \(v\), then \((a,H)\) has \(Z\)-coordinate
\[
    \bigl(u,\;q_{\mathbf{Beh}}(\partial_aH)\bigr),
\]
not
\[
    \bigl(v,\;q_{\mathbf{Beh}}(H)\bigr).
\]
This convention is used whenever derivative domains are involved in regular
image or factorization statements in \(\mathcal E/Z\).

\begin{definition}[Local coequational theory]
\label{def:local-coequational-theory}
    A \textbf{local coequational theory} over \((\mathbb A,\mathbb B)\) is a
    subrecognizer
    \[
        V\hookrightarrow \mathbf{Beh}^{\mathrm{can}}_{\mathbb A,\mathbb B}.
    \]
    Equivalently, it is a subobject
    \[
        j_V:V\hookrightarrow\mathbf{Beh}_0
    \]
    in \(\mathcal E/Z\) such that the canonical derivative map restricts to a
    transition
    \[
        \delta_V:
        A_1\times_{t_A,A_0,p_V}V\to V.
    \]
    The canonical output map then restricts to
    \[
        \omega_V:
        A_1\times_{t_A,A_0,p_V}V\to B_1.
    \]
\end{definition}

The derivative-closedness condition is the factorization
\[
\begin{tikzcd}[column sep=large, row sep=large]
    A_1\times_{t_A,A_0,p_V}V
        \arrow[rr, "{\partial(1\times j_V)}"]
        \arrow[dr, "\delta_V"']
    &&
    \mathbf{Beh}_0
    \\
    &
    V
        \arrow[ur, hook, "j_V"']
    &
\end{tikzcd}
\]
in \(\mathcal E/Z\).

When \(V\hookrightarrow\mathbf{Beh}_0\) is a local coequational theory, we
write
\[
    \delta_V(a,H)=\partial_aH
\]
and
\[
    \omega_V(a,H)=\omega_{\mathbf{Beh}}(a,H).
\]

\begin{definition}[Models]
\label{def:models-coequational-theory}
    Let
    \[
        V\hookrightarrow\mathbf{Beh}^{\mathrm{can}}
    \]
    be a local coequational theory. A recognizer \(P\) \textbf{models} \(V\) if
    its terminal semantic map
    \[
        \beta_P:P\to\mathbf{Beh}^{\mathrm{can}}
    \]
    factors through \(V\). We write
    \[
        \mathsf{Mod}(V)
    \]
    for the full subcategory of
    \[
        \mathsf{Rec}_{\mathbb A,\mathbb B}
    \]
    on the recognizers that model \(V\).
\end{definition}

\begin{lemma}[Model factorizations are strict]
\label{lem:model-factorizations-strict}
    If \(P\models V\), then the unique factorization
    \[
        \bar\beta_P:P\to V
    \]
    satisfying
    \[
        j_V\bar\beta_P=\beta_P
    \]
    is a strict homomorphism.
\end{lemma}

\begin{proof}
    \leavevmode
    \begin{enumerate}
        \item \textbf{Check transition preservation after inclusion.}
        Both
        \[
            j_V\bar\beta_P\delta_P
        \]
        and
        \[
            j_V\delta_V(1\times\bar\beta_P)
        \]
        are equal to the transition-preservation composite for the strict map
        \[
            \beta_P:P\to\mathbf{Beh}^{\mathrm{can}}.
        \]
        Since \(j_V\) is monic, the transition equation holds in \(V\).

        \item \textbf{Check output preservation.}
        The output equation for \(\bar\beta_P\) follows directly from the output
        equation for \(\beta_P\), since \(\omega_V\) is the restriction of
        \(\omega_{\mathbf{Beh}}\).

        \item \textbf{Check the map over \(Z\).}
        The equation
        \[
            j_V\bar\beta_P=\beta_P
        \]
        and the fact that both \(j_V\) and \(\beta_P\) are maps over \(Z\)
        imply that \(\bar\beta_P\) is a map over \(Z\).
    \end{enumerate}
\end{proof}

\subsection{Regular images and semantic images}
\label{subsec:regular-semantic-images}

\begin{proposition}[Regular images are created by the carrier slice]
\label{prop:regular-image-factorization-recognizers}
    Let
    \[
        f:P\to Q
    \]
    be a strict homomorphism of Mealy recognizers. Factor the underlying map in
    \(\mathcal E/Z\) as
    \[
        P\xrightarrow{e}I\xrightarrow{m}Q,
    \]
    where \(e\) is a regular epimorphism and \(m\) is a monomorphism. Then \(I\)
    carries a unique Mealy recognizer structure such that
    \[
        e:P\twoheadrightarrow I
    \]
    and
    \[
        m:I\hookrightarrow Q
    \]
    are strict homomorphisms.
\end{proposition}

\begin{proof}
    \leavevmode
    \begin{enumerate}
        \item \textbf{Pull back the image cover to transition domains.}
        Since regular epimorphisms are stable under pullback, the induced map
        \[
            A_1\times_{t_A,A_0,p_P}P
            \twoheadrightarrow
            A_1\times_{t_A,A_0,p_I}I
        \]
        is a regular epimorphism. It sends
        \[
            (a,x)\longmapsto(a,ex).
        \]

        \item \textbf{Descend the transition.}
        Define the transition on the cover by
        \[
            \delta_I(a,ex):=e(\delta_P(a,x)).
        \]
        This is independent of the chosen representative because, if \(ex=ex'\),
        then
        \[
        \begin{aligned}
            m e(\delta_P(a,x))
            &=
            f(\delta_P(a,x)) \\
            &=
            \delta_Q(a,fx) \\
            &=
            \delta_Q(a,mex) \\
            &=
            \delta_Q(a,mex') \\
            &=
            f(\delta_P(a,x')) \\
            &=
            m e(\delta_P(a,x')).
        \end{aligned}
        \]
        Since \(m\) is monic, the two values in \(I\) agree.

        \item \textbf{Descend the output.}
        Define the output on the same cover by
        \[
            \omega_I(a,ex):=\omega_P(a,x).
        \]
        If \(ex=ex'\), then \(fx=fx'\), and strictness of \(f\) gives
        \[
            \omega_P(a,x)=\omega_Q(a,fx)=\omega_Q(a,fx')=\omega_P(a,x').
        \]
        Hence the output descends.

        \item \textbf{Verify typing and laws.}
        The typing equations, unit law, and multiplication law hold after
        precomposition with the regular-epimorphic covers from \(P\). Hence they
        hold in \(I\).

        \item \textbf{Conclude creation of the factorization.}
        The defining equations make both
        \[
            e:P\to I
            \qquad\text{and}\qquad
            m:I\to Q
        \]
        strict homomorphisms, and uniqueness follows from epimorphy of \(e\)
        and monicity of \(m\).
    \end{enumerate}
\end{proof}

\begin{definition}[Semantic image]
\label{def:semantic-image}
    For a recognizer \(P\), its \textbf{semantic image} is the regular image
    factorization of the unique map to the terminal recognizer:
    \[
        P\xrightarrow{e_P}\operatorname{Beh}(P)
        \xrightarrow{m_P}
        \mathbf{Beh}^{\mathrm{can}}_{\mathbb A,\mathbb B}.
    \]
    On carriers, this is the regular image factorization of
    \[
        \beta_P:P\to\mathbf{Beh}_0
    \]
    in \(\mathcal E/Z\).
\end{definition}

The semantic image is displayed by
\[
\begin{tikzcd}[column sep=large, row sep=large]
    P
        \arrow[rr, "\beta_P"]
        \arrow[dr, two heads, "e_P"']
    &&
    \mathbf{Beh}_0
    \\
    &
    \operatorname{Beh}(P)
        \arrow[ur, hook, "m_P"']
    &
\end{tikzcd}
\qquad
\text{in }\mathcal E/Z.
\]

\begin{corollary}[Semantic images are coequational theories]
\label{cor:semantic-images-are-coeq-theories}
    For every recognizer \(P\), the semantic image
    \[
        \operatorname{Beh}(P)\hookrightarrow\mathbf{Beh}^{\mathrm{can}}
    \]
    is a local coequational theory. Moreover,
    \[
        P\twoheadrightarrow\operatorname{Beh}(P)
    \]
    is a regular quotient recognizer.
\end{corollary}

\begin{proof}
    \leavevmode
    \begin{enumerate}
        \item \textbf{Apply image creation.}
        The semantic image is the regular image of the strict homomorphism
        \[
            \beta_P:P\to\mathbf{Beh}^{\mathrm{can}}.
        \]
        Proposition~\ref{prop:regular-image-factorization-recognizers} creates
        this image in the recognizer category.

        \item \textbf{Identify the image as a subrecognizer of the terminal object.}
        The monic part
        \[
            \operatorname{Beh}(P)\hookrightarrow\mathbf{Beh}^{\mathrm{can}}
        \]
        is a subrecognizer, hence a local coequational theory.

        \item \textbf{Identify the quotient map.}
        The regular-epimorphic part
        \[
            P\twoheadrightarrow\operatorname{Beh}(P)
        \]
        is strict by the same proposition, hence is a regular quotient
        recognizer.
    \end{enumerate}
\end{proof}

\begin{lemma}[Meets of local theories]
\label{lem:meets-local-theories}
    If
    \[
        S,T\hookrightarrow\mathbf{Beh}^{\mathrm{can}}
    \]
    are local coequational theories, then their meet
    \[
        S\wedge T\hookrightarrow\mathbf{Beh}^{\mathrm{can}}
    \]
    is a local coequational theory and is the pullback subrecognizer.
\end{lemma}

\begin{proof}
    \leavevmode
    \begin{enumerate}
        \item \textbf{Take the pullback of subrecognizers.}
        The meet is the pullback in the carrier slice:
        \[
        \begin{tikzcd}[column sep=large, row sep=large]
            S\wedge T
                \arrow[r]
                \arrow[d]
                \arrow[dr, phantom, "\lrcorner", very near start]
            &
            T
                \arrow[d, hook]
            \\
            S
                \arrow[r, hook]
            &
            \mathbf{Beh}^{\mathrm{can}} .
        \end{tikzcd}
        \]

        \item \textbf{Restrict the recognizer structure.}
        Since both \(S\) and \(T\) are subrecognizers of
        \(\mathbf{Beh}^{\mathrm{can}}\), the transition and output of the
        canonical recognizer restrict to the intersection. The laws hold by
        monicity.

        \item \textbf{Conclude the universal property.}
        The resulting recognizer is the pullback in
        \(\mathsf{Rec}_{\mathbb A,\mathbb B}\), and its carrier is the meet in
        \(\mathcal E/Z\).
    \end{enumerate}
\end{proof}

\subsection{Small coproducts of recognizers}
\label{subsec:small-coproducts-recognizers}

\begin{proposition}[Small coproducts are created by the carrier slice]
\label{prop:small-coproducts-recognizers}
    A small family
    \[
        (P_i)_{i\in I}
    \]
    of recognizers over \((\mathbb A,\mathbb B)\) has a coproduct in
    \[
        \mathsf{Rec}_{\mathbb A,\mathbb B}.
    \]
    Its carrier is the coproduct
    \[
        \coprod_{i\in I}P_i
    \]
    in the slice
    \[
        \mathcal E/Z,
    \]
    and its transition and output maps are induced componentwise.
\end{proposition}

\begin{proof}
    \leavevmode
    \begin{enumerate}
        \item \textbf{Use preservation of coproducts by pullback.}
        In a Grothendieck topos, pullback preserves small coproducts. Hence
        \[
            A_1\times_{t_A,A_0,p_{\coprod_iP_i}}
            \coprod_iP_i
            \cong
            \coprod_i
            \left(
                A_1\times_{t_A,A_0,p_{P_i}}P_i
            \right).
        \]

        \item \textbf{Define transition and output componentwise.}
        The transition and output maps are induced by the component maps
        \[
            \delta_{P_i},
            \qquad
            \omega_{P_i}.
        \]

        \item \textbf{Verify the laws.}
        The Mealy laws hold after restriction to each coproduct summand. Since
        the summand inclusions are jointly epic, the laws hold on the coproduct.

        \item \textbf{Check the universal property.}
        A strict homomorphism
        \[
            \coprod_iP_i\to Q
        \]
        is equivalently a jointly compatible family of strict homomorphisms
        \[
            P_i\to Q.
        \]
        This is the universal property of the coproduct in \(\mathcal E/Z\)
        together with the componentwise recognizer structure.
    \end{enumerate}
\end{proof}

\subsection{Local closure}
\label{subsec:local-coequational-closure}

\begin{theorem}[Local coequational closure]
\label{thm:local-coequational-closure}
    Let
    \[
        V\hookrightarrow\mathbf{Beh}^{\mathrm{can}}
    \]
    be a local coequational theory. Then:
    \begin{enumerate}
        \item \(\mathsf{Mod}(V)\) is closed under strict homomorphic preimages.
        \item \(\mathsf{Mod}(V)\) is closed under regular quotient recognizers.
        \item \(\mathsf{Mod}(V)\) is closed under small coproducts.
        \item \(V\), regarded as a restricted canonical recognizer, is itself a
        model of \(V\), and
        \[
            \operatorname{Beh}(V)=V.
        \]
    \end{enumerate}
    Consequently the assignment
    \[
        V\longmapsto\mathsf{Mod}(V)
    \]
    is monotone and order-reflecting.
\end{theorem}

\begin{proof}
    Let
    \[
        j_V:V\hookrightarrow\mathbf{Beh}^{\mathrm{can}}
    \]
    denote the inclusion.

    \begin{enumerate}
        \item \textbf{Closure under strict homomorphic preimages.}
        Let
        \[
            f:P\to Q
        \]
        be strict and suppose \(Q\models V\). Choose
        \[
            \bar\beta_Q:Q\to V
        \]
        with
        \[
            j_V\bar\beta_Q=\beta_Q.
        \]
        Then
        \[
            \beta_P=\beta_Qf=j_V\bar\beta_Qf,
        \]
        so \(P\models V\).

        \item \textbf{Closure under regular quotient recognizers.}
        Let
        \[
            e:P\twoheadrightarrow Q
        \]
        be a regular quotient recognizer and suppose \(P\models V\). Choose
        \[
            \bar\beta_P:P\to V
        \]
        with
        \[
            j_V\bar\beta_P=\beta_P.
        \]
        Since
        \[
            \beta_P=\beta_Qe,
        \]
        the map \(\bar\beta_P\) is constant on the kernel pair of \(e\) after
        composing with the monomorphism \(j_V\). Hence it descends uniquely to
        \[
            \bar\beta_Q:Q\to V
        \]
        with
        \[
            \bar\beta_Qe=\bar\beta_P.
        \]
        Since \(e\) is epic,
        \[
            j_V\bar\beta_Q=\beta_Q,
        \]
        so \(Q\models V\).

        \item \textbf{Closure under small coproducts.}
        Let
        \[
            (P_i)_{i\in I}
        \]
        be a small family of models. For each \(i\), the factorization
        \[
            \bar\beta_i:P_i\to V
        \]
        is strict by Lemma~\ref{lem:model-factorizations-strict}. By the
        coproduct universal property in
        \(\mathsf{Rec}_{\mathbb A,\mathbb B}\), these strict maps induce a
        strict map
        \[
            \coprod_iP_i\to V.
        \]
        Therefore the coproduct models \(V\).

        \item \textbf{The theory models itself.}
        The identity map
        \[
            V\to V
        \]
        exhibits \(V\) as a model of itself. By
        Lemma~\ref{lem:canonical-behaviour-map-identity}, the semantic map of a
        restricted canonical subrecognizer is its inclusion
        \[
            j_V:V\hookrightarrow\mathbf{Beh}^{\mathrm{can}}.
        \]
        Hence
        \[
            \operatorname{Beh}(V)=V.
        \]

        \item \textbf{Prove monotonicity and order-reflection.}
        If \(V\leq W\), then every factorization through \(V\) is a
        factorization through \(W\), so
        \[
            \mathsf{Mod}(V)\subseteq\mathsf{Mod}(W).
        \]
        Conversely, if
        \[
            \mathsf{Mod}(V)\subseteq\mathsf{Mod}(W),
        \]
        then \(V\in\mathsf{Mod}(V)\), hence \(V\in\mathsf{Mod}(W)\). Therefore
        the inclusion
        \[
            V\hookrightarrow\mathbf{Beh}^{\mathrm{can}}
        \]
        factors through
        \[
            W\hookrightarrow\mathbf{Beh}^{\mathrm{can}},
        \]
        so \(V\leq W\).
    \end{enumerate}
\end{proof}

\subsection{The local coequational Birkhoff theorem}
\label{subsec:local-coequational-birkhoff}

Let
\[
    \mathcal C\subseteq\mathsf{Rec}_{\mathbb A,\mathbb B}
\]
be a replete full subcategory. Define
\[
    \mathcal S_{\mathcal C}
    :=
    \{
        S\hookrightarrow\mathbf{Beh}^{\mathrm{can}}
        \mid
        S=\operatorname{Beh}(P)
        \text{ for some }P\in\mathcal C
    \}.
\]
This is a set of subobjects of \(\mathbf{Beh}^{\mathrm{can}}\) in
\(\mathcal E/Z\). Define
\[
    V_{\mathcal C}
    :=
    \bigvee_{S\in\mathcal S_{\mathcal C}}S
    \hookrightarrow\mathbf{Beh}^{\mathrm{can}}.
\]

\begin{lemma}[Joins of local coequational theories]
\label{lem:joins-derivative-closed}
    Let
    \[
        (S_i\hookrightarrow\mathbf{Beh}^{\mathrm{can}})_{i\in I}
    \]
    be a small family of local coequational theories. Then their join
    \[
        S:=\bigvee_{i\in I}S_i
    \]
    in
    \[
        \operatorname{Sub}_{\mathcal E/Z}(\mathbf{Beh}_0)
    \]
    is again a local coequational theory.
\end{lemma}

\begin{proof}
    \leavevmode
    \begin{enumerate}
        \item \textbf{Represent the join.}
        In the Grothendieck topos \(\mathcal E/Z\), the join is represented by
        the regular image of the coproduct map
        \[
            \coprod_iS_i\to\mathbf{Beh}_0.
        \]
        Write this image factorization as
        \[
        \begin{tikzcd}[column sep=large, row sep=large]
            \coprod_iS_i
                \arrow[rr]
                \arrow[dr, two heads]
            &&
            \mathbf{Beh}_0
            \\
            &
            S
                \arrow[ur, hook]
            &
        \end{tikzcd}
        \]
        in \(\mathcal E/Z\).

        \item \textbf{Regard the coproduct as a recognizer.}
        By Proposition~\ref{prop:small-coproducts-recognizers}, the coproduct
        \[
            \coprod_iS_i
        \]
        carries the componentwise recognizer structure. Since each
        \[
            S_i\hookrightarrow\mathbf{Beh}^{\mathrm{can}}
        \]
        is strict, the induced coproduct map to
        \(\mathbf{Beh}^{\mathrm{can}}\) is strict.

        \item \textbf{Take the regular image in recognizers.}
        By Proposition~\ref{prop:regular-image-factorization-recognizers}, the
        regular image of this strict map is created in \(\mathcal E/Z\) and
        carries a recognizer structure. Hence
        \[
            S\hookrightarrow\mathbf{Beh}^{\mathrm{can}}
        \]
        is a subrecognizer.

        \item \textbf{Conclude local coequationality.}
        Being a subrecognizer of the canonical behaviour recognizer is exactly
        being a local coequational theory.
    \end{enumerate}
\end{proof}

\begin{theorem}[Local coequational Birkhoff theorem]
\label{thm:local-coequational-birkhoff}
    Let
    \[
        \mathcal C\subseteq\mathsf{Rec}_{\mathbb A,\mathbb B}
    \]
    be a replete full subcategory. Then \(\mathcal C\) is of the form
    \[
        \mathcal C=\mathsf{Mod}(V)
    \]
    for a local coequational theory
    \[
        V\hookrightarrow\mathbf{Beh}^{\mathrm{can}}_{\mathbb A,\mathbb B}
    \]
    if and only if \(\mathcal C\) is closed under:
    \begin{enumerate}
        \item strict homomorphic preimages;
        \item regular quotient recognizers;
        \item small coproducts.
    \end{enumerate}
    When these conditions hold, the defining theory is
    \[
        V_{\mathcal C}
        =
        \bigvee_{P\in\mathcal C}\operatorname{Beh}(P)
        \hookrightarrow\mathbf{Beh}^{\mathrm{can}}.
    \]
\end{theorem}

\begin{proof}
    \leavevmode
    \begin{enumerate}
        \item \textbf{Prove necessity.}
        If
        \[
            \mathcal C=\mathsf{Mod}(V),
        \]
        then the desired closure properties are
        Theorem~\ref{thm:local-coequational-closure}.

        \item \textbf{Construct the candidate theory.}
        Conversely, suppose \(\mathcal C\) is closed under strict homomorphic
        preimages, regular quotient recognizers, and small coproducts. Define
        \[
            V_{\mathcal C}
            :=
            \bigvee_{P\in\mathcal C}\operatorname{Beh}(P)
            \hookrightarrow\mathbf{Beh}^{\mathrm{can}}.
        \]
        Each \(\operatorname{Beh}(P)\) is a local coequational theory by
        Corollary~\ref{cor:semantic-images-are-coeq-theories}, so
        \(V_{\mathcal C}\) is a local coequational theory by
        Lemma~\ref{lem:joins-derivative-closed}.

        \item \textbf{Show \(\mathcal C\subseteq\mathsf{Mod}(V_{\mathcal C})\).}
        If \(P\in\mathcal C\), then
        \[
            \operatorname{Beh}(P)\leq V_{\mathcal C}.
        \]
        Since \(\beta_P\) factors through \(\operatorname{Beh}(P)\), it factors
        through \(V_{\mathcal C}\). Hence \(P\models V_{\mathcal C}\).

        \item \textbf{Let a model be given.}
        Let
        \[
            X\in\mathsf{Mod}(V_{\mathcal C}).
        \]
        Then
        \[
            \operatorname{Beh}(X)\leq V_{\mathcal C}.
        \]

        \item \textbf{Decompose the semantic image by intersections.}
        Since the subobject lattice of \(\mathcal E/Z\) is a complete Heyting
        algebra,
        \[
        \begin{aligned}
            \operatorname{Beh}(X)
            &=
            \operatorname{Beh}(X)\wedge V_{\mathcal C} \\
            &=
            \operatorname{Beh}(X)\wedge
            \bigvee_{S\in\mathcal S_{\mathcal C}}S \\
            &=
            \bigvee_{S\in\mathcal S_{\mathcal C}}
            \bigl(\operatorname{Beh}(X)\wedge S\bigr).
        \end{aligned}
        \]

        \item \textbf{Show each intersection belongs to \(\mathcal C\).}
        Choose \(P_S\in\mathcal C\) with
        \[
            S=\operatorname{Beh}(P_S).
        \]
        Since
        \[
            P_S\twoheadrightarrow S
        \]
        is a regular quotient recognizer, closure under regular quotients gives
        \[
            S\in\mathcal C.
        \]
        By Lemma~\ref{lem:meets-local-theories},
        \[
            \operatorname{Beh}(X)\wedge S
        \]
        is a subrecognizer of \(S\). Closure under strict homomorphic preimages,
        applied to the monic strict map
        \[
            \operatorname{Beh}(X)\wedge S\hookrightarrow S,
        \]
        gives
        \[
            \operatorname{Beh}(X)\wedge S\in\mathcal C.
        \]

        \item \textbf{Use coproducts and regular images.}
        Form
        \[
            C_X:=
            \coprod_{S\in\mathcal S_{\mathcal C}}
            \bigl(\operatorname{Beh}(X)\wedge S\bigr).
        \]
        By closure under small coproducts,
        \[
            C_X\in\mathcal C.
        \]
        The regular image of the canonical strict map
        \[
            C_X\to\mathbf{Beh}^{\mathrm{can}}
        \]
        is precisely
        \[
            \operatorname{Beh}(X),
        \]
        by the join calculation. Hence
        \[
            C_X\twoheadrightarrow\operatorname{Beh}(X)
        \]
        is a regular quotient recognizer, so
        \[
            \operatorname{Beh}(X)\in\mathcal C.
        \]

        \item \textbf{Pull back along the semantic quotient.}
        The semantic image factorization gives a regular quotient recognizer
        \[
            X\twoheadrightarrow\operatorname{Beh}(X).
        \]
        Since \(\operatorname{Beh}(X)\in\mathcal C\) and \(\mathcal C\) is
        closed under strict homomorphic preimages, it follows that
        \[
            X\in\mathcal C.
        \]

        \item \textbf{Conclude equality.}
        Thus
        \[
            \mathsf{Mod}(V_{\mathcal C})\subseteq\mathcal C.
        \]
        Together with the opposite inclusion, this gives
        \[
            \mathcal C=\mathsf{Mod}(V_{\mathcal C}).
        \]
    \end{enumerate}
\end{proof}

\subsection{Single-object and stateless specializations}
\label{subsec:coeq-single-object-stateless}

Suppose first that \(\mathbb A\) and \(\mathbb B\) are one-object internal
categories corresponding to monoid objects \(M\) and \(N\). A recognizer is a
right \(M\)-action object \(P\), written
\[
    x\cdot m:=\delta_P(m,x),
\]
together with an output cocycle
\[
    \omega_P:M\times P\to N
\]
satisfying
\[
    x\cdot 1=x,
    \qquad
    x\cdot(mn)=(x\cdot m)\cdot n,
\]
and
\[
    \omega_P(1,x)=1,
    \qquad
    \omega_P(mn,x)
    =
    \omega_P(m,x)\,\omega_P(n,x\cdot m).
\]
A local coequational theory is a subobject of represented residual behaviours
closed under all residual derivatives. With the convention
\[
    b\circ a=ba,
\]
derivatives are placed on the left:
\[
    (\partial_mH)(r)=H(mr).
\]

In the stateless, or map-representable, case a recognizer is induced by an
internal functor
\[
    F:\mathbb A\to\mathbb B.
\]
Its carrier is
\[
    A_0\xleftarrow{1}A_0\xrightarrow{F_0}B_0,
\]
with
\[
    \delta_F(a,v)=s_A(a),
    \qquad
    \omega_F(a,v)=F_1(a).
\]
The residual semantics of \(v\in A_0\) is
\[
    H^F_v:\mathbb A/v\to\mathbb B
\]
defined by
\[
    H^F_v(r:u\to v)=F_0(u)
\]
and, on a residual arrow
\[
    w\xrightarrow{a}u\xrightarrow{r}v,
\]
by
\[
    H^F_v(a,r)=F_1(a):F_0(w)\to F_0(u).
\]
Thus a stateless recognizer models \(V\) precisely when all residual behaviours
\[
    H^F_v
\]
lie in \(V\).

\section{Structured output and algebraic closure}
\label{sec:structured-output-algebraic-closure}

This section explains how additional closure operations on behaviours arise
from algebraic structure on the output category.

\subsection{Structured output categories}
\label{subsec:structured-output-categories}

\begin{definition}[Structured output category]
\label{def:structured-output-category}
    Let \(\mathbb T\) be a single-sorted finite-product theory. A
    \textbf{\(\mathbb T\)-structured output category} is a \(\mathbb T\)-algebra
    object in internal categories, equivalently a finite-product-preserving
    functor
    \[
        \mathbb T\longrightarrow \mathbf{Cat}_{\mathrm{int}}(\mathcal E).
    \]
    We write
    \[
        \mathbb B
    \]
    for the internal category assigned to the generic object of \(\mathbb T\).
    Thus each operation of arity \(n\) in \(\mathbb T\) determines an internal
    functor
    \[
        \theta:\mathbb B^n\to\mathbb B,
    \]
    and the equations of \(\mathbb T\) hold internally.
\end{definition}

\subsection{Pointwise operations on behaviours}
\label{subsec:pointwise-operations-behaviours}

For \(n>0\), put
\[
    \mathbf{Beh}_0^{(n)}
    :=
    \underbrace{
    \mathbf{Beh}_0\times_{A_0}\cdots\times_{A_0}\mathbf{Beh}_0
    }_{n}.
\]
A generalized element is a tuple
\[
    (H_1,\dots,H_n)
\]
of residual behaviours
\[
    H_i:\mathbb A/v\to\mathbb B
\]
over the same input object \(v\). For \(n=0\), set
\[
    \mathbf{Beh}_0^{(0)}:=A_0.
\]

\begin{definition}[Pointwise structured behaviour]
\label{def:pointwise-structured-behaviour}
    Let
    \[
        \theta:\mathbb B^n\to\mathbb B
    \]
    be an operation of \(\mathbb T\). The induced map
    \[
        \theta_{\mathbf{Beh}}:
        \mathbf{Beh}_0^{(n)}\to\mathbf{Beh}_0
    \]
    is defined by
    \[
        \theta_{\mathbf{Beh}}(H_1,\dots,H_n)
        =
        \theta\circ\langle H_1,\dots,H_n\rangle
    \]
    for \(n>0\). For \(n=0\), the nullary operation
    \[
        \theta:1\to\mathbb B
    \]
    induces the constant residual behaviour
    \[
        \mathbb A/v\to 1\xrightarrow{\theta}\mathbb B
    \]
    over each \(v\in A_0\).
\end{definition}

The construction is summarized by
\[
\begin{tikzcd}[column sep=huge, row sep=large]
    \mathbb A/v
        \arrow[r, "{\langle H_1,\dots,H_n\rangle}"]
        \arrow[rr, "{\theta_{\mathbf{Beh}}(H_i)_i}"', bend right=14]
    &
    \mathbb B^n
        \arrow[r, "\theta"]
    &
    \mathbb B .
\end{tikzcd}
\]

The map
\[
    \theta_{\mathbf{Beh}}
\]
is naturally a map over \(A_0\). It is not usually a map over
\[
    Z=A_0\times B_0.
\]
For \(n>0\), the \(B_0\)-coordinate of
\[
    \theta_{\mathbf{Beh}}(H_1,\dots,H_n)
\]
is
\[
    \theta_0(q_{\mathbf{Beh}}(H_1),\dots,q_{\mathbf{Beh}}(H_n)).
\]
For \(n=0\), it is the object of \(\mathbb B\) selected by the nullary
operation.

\begin{lemma}[Structured operations commute with residual semantics]
\label{lem:structured-operations-semantic}
    For every operation
    \[
        \theta:\mathbb B^n\to\mathbb B,
    \]
    there is an induced operation on recognizers
    \[
        (P_1,\dots,P_n)
        \longmapsto
        \theta(P_1,\dots,P_n)
    \]
    and the residual semantics is compatible with this operation:
    \[
        \beta_{\theta(P_i)_i}
        =
        \theta_{\mathbf{Beh}}(\beta_{P_i})_i.
    \]
\end{lemma}

\begin{proof}
    \leavevmode
    \begin{enumerate}
        \item \textbf{Construct the carrier.}
        For \(n>0\), the carrier is
        \[
            P_\theta:=P_1\times_{A_0}\cdots\times_{A_0}P_n.
        \]
        Its output projection is
        \[
            q_\theta(x_1,\dots,x_n)
            =
            \theta_0(q_1x_1,\dots,q_nx_n).
        \]
        For \(n=0\), the carrier is \(A_0\), with output projection the
        constant object selected by the nullary operation.

        \item \textbf{Define transition and output.}
        For \(n>0\),
        \[
            \delta_\theta(a,(x_i)_i)
            =
            (\delta_i(a,x_i))_i
        \]
        and
        \[
            \omega_\theta(a,(x_i)_i)
            =
            \theta_1(\omega_i(a,x_i))_i.
        \]
        For \(n=0\), the transition sends \(v\) to \(u\) along
        \(a:u\to v\), and the output is the identity arrow selected by the
        nullary functor.

        \item \textbf{Verify the Mealy laws structurally.}
        The transition laws are inherited componentwise. The output laws follow
        from functoriality of
        \[
            \theta:\mathbb B^n\to\mathbb B.
        \]

        \item \textbf{Compare residual behaviours.}
        For a residual object
        \[
            r:u\to v,
        \]
        one has
        \[
        \begin{aligned}
            q_\theta(\delta_\theta(r,(x_i)_i))
            &=
            \theta_0(q_i(\delta_i(r,x_i)))_i,
        \end{aligned}
        \]
        which is the object value of
        \[
            \theta_{\mathbf{Beh}}(\beta_{P_i}(x_i))_i.
        \]
        The same calculation on residual arrows gives
        \[
            \omega_\theta(a,\delta_\theta(r,(x_i)_i))
            =
            \theta_1(\omega_i(a,\delta_i(r,x_i)))_i.
        \]
        Therefore
        \[
            \beta_{\theta(P_i)_i}
            =
            \theta_{\mathbf{Beh}}(\beta_{P_i})_i.
        \]
    \end{enumerate}
\end{proof}

\begin{lemma}[Structured operations preserve derivatives and canonical output]
\label{lem:operations-preserve-derivatives-output}
    For every operation
    \[
        \theta:\mathbb B^n\to\mathbb B
    \]
    and every arrow
    \[
        a:u\to v
    \]
    of \(\mathbb A\),
    \[
        \partial_a\bigl(\theta_{\mathbf{Beh}}(H_i)_i\bigr)
        =
        \theta_{\mathbf{Beh}}(\partial_aH_i)_i
    \]
    and
    \[
        \omega_{\mathbf{Beh}}
        \bigl(a,\theta_{\mathbf{Beh}}(H_i)_i\bigr)
        =
        \theta_1\bigl(\omega_{\mathbf{Beh}}(a,H_i)\bigr)_i.
    \]
\end{lemma}

\begin{proof}
    \leavevmode
    \begin{enumerate}
        \item \textbf{Use precomposition for derivatives.}
        Since
        \[
            \partial_aH=H\circ a_*,
        \]
        one has
        \[
        \begin{aligned}
            \partial_a\bigl(\theta\circ\langle H_i\rangle_i\bigr)
            &=
            \bigl(\theta\circ\langle H_i\rangle_i\bigr)\circ a_* \\
            &=
            \theta\circ\langle H_i\circ a_*\rangle_i \\
            &=
            \theta_{\mathbf{Beh}}(\partial_aH_i)_i.
        \end{aligned}
        \]

        \item \textbf{Use functoriality for canonical output.}
        The canonical output is the value on the residual morphism
        \[
            a\to 1_v.
        \]
        Applying
        \[
            \theta\circ\langle H_i\rangle_i
        \]
        gives
        \[
            \theta_1(H_i(a,1_v))_i
            =
            \theta_1(\omega_{\mathbf{Beh}}(a,H_i))_i.
        \]

        \item \textbf{Handle nullary operations.}
        In the nullary case both claims say that a constant functor remains
        constant under reindexing and sends residual morphisms to identity
        arrows selected by the nullary functor.
    \end{enumerate}
\end{proof}

\subsection{\texorpdfstring{\(\mathbb T\)}{T}-coequational theories}
\label{subsec:T-coequational-theories}

\begin{definition}[\(\mathbb T\)-coequational theory]
\label{def:T-coequational-theory}
    Let \(\mathbb B\) be a \(\mathbb T\)-structured output category. A
    \textbf{local \(\mathbb T\)-coequational theory} is a local coequational
    theory
    \[
        V\hookrightarrow\mathbf{Beh}^{\mathrm{can}}_{\mathbb A,\mathbb B}
    \]
    whose underlying subobject is closed under the pointwise
    \(\mathbb T\)-algebra structure.

    Explicitly, for an \(n\)-ary operation
    \[
        \theta:\mathbb B^n\to\mathbb B,
    \]
    write
    \[
        V^{(n)}_\theta
        :=
        \underbrace{V\times_{A_0}\cdots\times_{A_0}V}_{n}
    \]
    for \(n>0\), equipped with the \(Z\)-structure
    \[
        (H_1,\dots,H_n)
        \longmapsto
        \Bigl(
            p_V(H_1),
            \theta_0(q_V(H_1),\dots,q_V(H_n))
        \Bigr).
    \]
    For \(n=0\), put
    \[
        V^{(0)}_\theta:=A_0
    \]
    with \(Z\)-coordinate
    \[
        v\longmapsto (v,\theta_0(*)).
    \]
    Closure under \(\theta\) means that
    \[
        \theta_{\mathbf{Beh}}:V^{(n)}_\theta\to\mathbf{Beh}_0
    \]
    factors through
    \[
        V\hookrightarrow\mathbf{Beh}_0
    \]
    as a map over \(Z\). The induced map is denoted
    \[
        \theta_V:V^{(n)}_\theta\to V.
    \]
\end{definition}

The closure condition is displayed by
\[
\begin{tikzcd}[column sep=large, row sep=large]
    V^{(n)}_\theta
        \arrow[rr, "{\theta_{\mathbf{Beh}}}"]
        \arrow[dr, dashed, "\theta_V"']
    &&
    \mathbf{Beh}_0
    \\
    &
    V
        \arrow[ur, hook, "j_V"']
    &
\end{tikzcd}
\qquad
\text{in }\mathcal E/Z.
\]

\subsection{\texorpdfstring{\(\mathbb T\)}{T}-products of recognizers}
\label{subsec:T-products-recognizers}

\begin{definition}[\(\mathbb T\)-product recognizer]
\label{def:T-product-recognizer}
    For recognizers
    \[
        P_1,\dots,P_n
    \]
    and an operation
    \[
        \theta:\mathbb B^n\to\mathbb B,
    \]
    the recognizer
    \[
        \theta(P_1,\dots,P_n)
    \]
    is the recognizer constructed in
    Lemma~\ref{lem:structured-operations-semantic}. For \(n=0\), it is the
    nullary recognizer with carrier \(A_0\) and constant output selected by
    \(\theta\).
\end{definition}

\begin{theorem}[\(\mathbb T\)-coequational closure]
\label{thm:T-coequational-closure}
    Let
    \[
        V\hookrightarrow\mathbf{Beh}^{\mathrm{can}}
    \]
    be a local \(\mathbb T\)-coequational theory. Then
    \[
        \mathsf{Mod}(V)
    \]
    is closed under strict homomorphic preimages, regular quotient recognizers,
    small coproducts, and \(\mathbb T\)-products.
\end{theorem}

\begin{proof}
    \leavevmode
    \begin{enumerate}
        \item \textbf{Use local closure.}
        Closure under strict homomorphic preimages, regular quotient
        recognizers, and small coproducts is
        Theorem~\ref{thm:local-coequational-closure}.

        \item \textbf{Use semantic compatibility.}
        Let
        \[
            P_1,\dots,P_n\in\mathsf{Mod}(V).
        \]
        Choose strict factorizations
        \[
            \bar\beta_{P_i}:P_i\to V
        \]
        with
        \[
            j_V\bar\beta_{P_i}=\beta_{P_i}.
        \]
        These induce a map from the carrier of
        \[
            \theta(P_1,\dots,P_n)
        \]
        to \(V^{(n)}_\theta\).

        \item \textbf{Use algebraic closure of \(V\).}
        Since \(V\) is closed under \(\theta\), there is a map
        \[
            \theta_V:V^{(n)}_\theta\to V
        \]
        with
        \[
            j_V\theta_V
            =
            \theta_{\mathbf{Beh}}\big|_{V^{(n)}_\theta}.
        \]
        By Lemma~\ref{lem:structured-operations-semantic},
        \[
            \beta_{\theta(P_i)_i}
            =
            \theta_{\mathbf{Beh}}(\beta_{P_i})_i.
        \]
        Hence
        \[
            \beta_{\theta(P_i)_i}
        \]
        factors through \(V\), so
        \[
            \theta(P_1,\dots,P_n)\in\mathsf{Mod}(V).
        \]

        \item \textbf{Handle nullary operations.}
        The nullary case is the same argument using
        \[
            V^{(0)}_\theta=A_0
        \]
        and the closure map
        \[
            A_0\to V.
        \]
    \end{enumerate}
\end{proof}

\begin{theorem}[\(\mathbb T\)-coequational Birkhoff theorem]
\label{thm:T-coequational-birkhoff}
    Let \(\mathbb B\) be a \(\mathbb T\)-structured output category, and let
    \[
        \mathcal C\subseteq\mathsf{Rec}_{\mathbb A,\mathbb B}
    \]
    be a replete full subcategory. Then \(\mathcal C\) is of the form
    \[
        \mathcal C=\mathsf{Mod}(V)
    \]
    for a local \(\mathbb T\)-coequational theory
    \[
        V\hookrightarrow\mathbf{Beh}^{\mathrm{can}}_{\mathbb A,\mathbb B}
    \]
    if and only if \(\mathcal C\) is closed under:
    \begin{enumerate}
        \item strict homomorphic preimages;
        \item regular quotient recognizers;
        \item small coproducts;
        \item \(\mathbb T\)-products.
    \end{enumerate}
\end{theorem}

\begin{proof}
    \leavevmode
    \begin{enumerate}
        \item \textbf{Prove necessity.}
        If
        \[
            \mathcal C=\mathsf{Mod}(V)
        \]
        for a local \(\mathbb T\)-coequational theory \(V\), then the closure
        properties are exactly Theorem~\ref{thm:T-coequational-closure}.

        \item \textbf{Construct the local candidate.}
        Conversely, assume \(\mathcal C\) has the four closure properties. By
        Theorem~\ref{thm:local-coequational-birkhoff},
        \[
            V_{\mathcal C}
            =
            \bigvee_{P\in\mathcal C}\operatorname{Beh}(P)
        \]
        is a local coequational theory and
        \[
            \mathcal C=\mathsf{Mod}(V_{\mathcal C}).
        \]
        It remains to prove that \(V_{\mathcal C}\) is closed under the
        pointwise \(\mathbb T\)-operations.

        \item \textbf{Cover the operation domain.}
        For an \(n\)-ary operation \(\theta\), the object
        \[
            (V_{\mathcal C})^{(n)}_\theta
        \]
        is formed over \(A_0\), with \(Z\)-coordinate induced by \(\theta_0\).
        Since finite meets distribute over small joins in a topos, it is
        covered by the coproduct, over tuples
        \[
            (P_1,\dots,P_n)
        \]
        of objects of \(\mathcal C\), of the objects
        \[
            \operatorname{Beh}(P_1)
            \times_{A_0}\cdots\times_{A_0}
            \operatorname{Beh}(P_n),
        \]
        equipped with the same \(\theta\)-induced \(Z\)-coordinate.

        \item \textbf{Compute on each cover component.}
        For each tuple \((P_i)_i\), the product of semantic quotient maps
        \[
            P_i\twoheadrightarrow\operatorname{Beh}(P_i)
        \]
        gives a regular-epimorphic cover of the corresponding product of
        semantic images. By Lemma~\ref{lem:structured-operations-semantic}, the
        pointwise operation sends this product into
        \[
            \operatorname{Beh}(\theta(P_1,\dots,P_n)).
        \]
        Since \(\mathcal C\) is closed under \(\mathbb T\)-products,
        \[
            \theta(P_1,\dots,P_n)\in\mathcal C,
        \]
        so
        \[
            \operatorname{Beh}(\theta(P_1,\dots,P_n))
            \leq
            V_{\mathcal C}.
        \]

        \item \textbf{Descend the factorization.}
        The operation map
        \[
            \theta_{\mathbf{Beh}}:
            (V_{\mathcal C})^{(n)}_\theta\to\mathbf{Beh}_0
        \]
        factors through \(V_{\mathcal C}\) after precomposition with a
        regular-epimorphic cover. By regular-epi descent for subobjects, it
        factors through \(V_{\mathcal C}\).

        \item \textbf{Handle nullary operations.}
        For \(n=0\), closure of \(\mathcal C\) under the nullary product
        recognizer implies that the corresponding constant residual behaviour
        has semantic image contained in \(V_{\mathcal C}\). Hence the nullary
        operation also factors through \(V_{\mathcal C}\).

        \item \textbf{Conclude.}
        Thus \(V_{\mathcal C}\) is a local \(\mathbb T\)-coequational theory,
        and
        \[
            \mathcal C=\mathsf{Mod}(V_{\mathcal C}).
        \]
    \end{enumerate}
\end{proof}

\subsection{Single-object and stateless specializations}
\label{subsec:structured-single-object-stateless}

In the one-object case, a \(\mathbb T\)-structured output category is a
\(\mathbb T\)-algebra object in one-object internal categories. Equivalently,
the output monoid \(N\) carries operations compatible with multiplication.
Pointwise operations on behaviours are computed objectwise on residual outputs
and arrowwise on residual output comparisons.

For one-object input and Boolean chaotic output
\[
    2_{\mathrm{ch}},
\]
local Boolean coequational theories are Boolean algebras of languages closed
under residual derivatives. They are not automatically theories of regular
languages; a regular-language specialization requires an additional
finite-state condition, or an explicit restriction to regular languages.

In the stateless case, a structured output operation
\[
    \theta:\mathbb B^n\to\mathbb B
\]
sends internal functors
\[
    F_i:\mathbb A\to\mathbb B
\]
to the composite functor
\[
    \mathbb A
    \xrightarrow{\langle F_i\rangle_i}
    \mathbb B^n
    \xrightarrow{\theta}
    \mathbb B.
\]
The corresponding stateless recognizer is precisely the
\(\mathbb T\)-product of the stateless recognizers induced by the \(F_i\).

\section{Global input reindexing}
\label{sec:global-input-reindexing}

This section explains how local coequational theories vary with the input
category. This is the categorical source of closure under inverse morphic
preimage in the classical one-object setting.

\subsection{Reindexing residual behaviours along input functors}
\label{subsec:input-reindexing-behaviours}

Fix \(\mathbb B\), and, when desired, a finite-product theory \(\mathbb T\) for
which \(\mathbb B\) is a \(\mathbb T\)-structured output category. Let
\(\mathsf{Inp}\) be a chosen category of internal input categories and internal
functors. For every
\[
    \mathbb A\in\mathsf{Inp},
\]
there is a canonical behaviour object
\[
    \mathbf{Beh}_{\mathbb A,\mathbb B,0}
\]
and a recognizer category
\[
    \mathsf{Rec}_{\mathbb A,\mathbb B}.
\]
We write
\[
    Z_{\mathbb A}:=A_0\times B_0.
\]

Let
\[
    F:\mathbb A'\to\mathbb A
\]
be an internal functor. For each \(v'\in A'_0\), there is an induced residual
slice functor
\[
    F_{/v'}:
    \mathbb A'/v'\to \mathbb A/F_0(v')
\]
sending
\[
    r':u'\to v'
\]
to
\[
    F_1(r'):F_0(u')\to F_0(v').
\]
On residual triangles it acts by applying \(F\):
\[
\begin{tikzcd}[column sep=large, row sep=large]
    w'
        \arrow[r, "a'"]
        \arrow[rr, "r'\circ a'", bend left=15]
    &
    u'
        \arrow[r, "r'"']
    &
    v'
\end{tikzcd}
\qquad\longmapsto\qquad
\begin{tikzcd}[column sep=large, row sep=large]
    F_0w'
        \arrow[r, "F_1a'"]
        \arrow[rr, "F_1(r'\circ a')"', bend right=15]
    &
    F_0u'
        \arrow[r, "F_1r'"]
    &
    F_0v'.
\end{tikzcd}
\]

By the representing property of the behaviour object, these residual slice
functors induce a reindexing map
\[
    F^\sharp:
    A'_0\times_{F_0,A_0,p_{\mathbf{Beh}}}
    \mathbf{Beh}_{\mathbb A,\mathbb B,0}
    \to
    \mathbf{Beh}_{\mathbb A',\mathbb B,0}
\]
given by
\[
    F^\sharp(v',H)=H\circ F_{/v'}.
\]
The domain is regarded over
\[
    Z_{\mathbb A'}=A'_0\times B_0
\]
by
\[
    (v',H)\longmapsto (v',q_{\mathbf{Beh}}(H)).
\]
With this convention, \(F^\sharp\) is a map over \(Z_{\mathbb A'}\).

\begin{lemma}[Reindexing preserves canonical structure]
\label{lem:global-reindexing-preserves-structure}
    The map \(F^\sharp\) preserves derivatives and canonical output
    comparisons. Explicitly,
    \[
        \partial_{a'}(F^\sharp(v',H))
        =
        F^\sharp(u',\partial_{F_1(a')}H)
    \]
    for every
    \[
        a':u'\to v',
    \]
    and
    \[
        \omega_{\mathbf{Beh}_{\mathbb A',\mathbb B}}
        (a',F^\sharp(v',H))
        =
        \omega_{\mathbf{Beh}_{\mathbb A,\mathbb B}}
        (F_1(a'),H).
    \]
    If \(\mathbb B\) is \(\mathbb T\)-structured, then \(F^\sharp\) also
    preserves the pointwise \(\mathbb T\)-structure.
\end{lemma}

\begin{proof}
    \leavevmode
    \begin{enumerate}
        \item \textbf{Compare residual transports.}
        The residual slice functors commute with postcomposition:
        \[
            F_{/v'}\circ a'_*
            =
            (F_1a')_*\circ F_{/u'}.
        \]
        Diagrammatically,
        \[
        \begin{tikzcd}[column sep=huge, row sep=large]
            \mathbb A'/u'
                \arrow[r, "a'_*"]
                \arrow[d, "F_{/u'}"']
            &
            \mathbb A'/v'
                \arrow[d, "F_{/v'}"]
            \\
            \mathbb A/F_0u'
                \arrow[r, "{(F_1a')_*}"']
            &
            \mathbb A/F_0v'.
        \end{tikzcd}
        \]
        Precomposition with this square gives the derivative identity.

        \item \textbf{Compare canonical outputs.}
        The canonical output of \(F^\sharp(v',H)\) along \(a'\) is the arrow
        assigned by
        \[
            H\circ F_{/v'}
        \]
        to the residual morphism
        \[
            a'\to 1_{v'}
        \]
        in \(\mathbb A'/v'\). Under \(F_{/v'}\), this residual morphism is sent
        to
        \[
            F_1(a')\to 1_{F_0(v')}
        \]
        in \(\mathbb A/F_0(v')\). Hence the output equality follows.

        \item \textbf{Compare pointwise operations.}
        The map \(F^\sharp\) acts by precomposition, while the pointwise
        \(\mathbb T\)-structure is defined by postcomposition with the
        structural functors
        \[
            \mathbb B^n\to\mathbb B.
        \]
        Therefore the two operations commute.
    \end{enumerate}
\end{proof}

\subsection{Input pullback of recognizers}
\label{subsec:input-pullback-recognizers}

\begin{definition}[Input pullback of recognizers]
\label{def:input-pullback-recognizer}
    Let
    \[
        P:\mathbb A\rightsquigarrow\mathbb B
    \]
    be a recognizer. Define
    \[
        F^*P:\mathbb A'\rightsquigarrow\mathbb B
    \]
    as follows. The carrier is
    \[
        F^*P
        :=
        A'_0\times_{F_0,A_0,p_P}P.
    \]
    A state is a pair
    \[
        (v',x)
    \]
    with
    \[
        F_0(v')=p_P(x).
    \]
    The input projection is
    \[
        p_{F^*P}(v',x)=v',
    \]
    and the output projection is
    \[
        q_{F^*P}(v',x)=q_P(x).
    \]
    For an arrow
    \[
        a':u'\to v',
    \]
    define
    \[
        \delta_{F^*P}(a',(v',x))
        =
        (u',\delta_P(F_1a',x)),
    \]
    and
    \[
        \omega_{F^*P}(a',(v',x))
        =
        \omega_P(F_1a',x).
    \]
\end{definition}

The carrier is the pullback
\[
\begin{tikzcd}[column sep=large, row sep=large]
    F^*P
        \arrow[r]
        \arrow[d]
        \arrow[dr, phantom, "\lrcorner", very near start]
    &
    P
        \arrow[d, "p_P"]
    \\
    A'_0
        \arrow[r, "F_0"']
    &
    A_0 .
\end{tikzcd}
\]

\begin{proposition}[Input pullback and semantics]
\label{prop:input-pullback-recognizer}
    The data above define a recognizer
    \[
        F^*P:\mathbb A'\rightsquigarrow\mathbb B.
    \]
    Moreover,
    \[
        \beta_{F^*P}
        =
        F^\sharp\beta_P
    \]
    in the sense that
    \[
        \beta_{F^*P}(v',x)
        =
        F^\sharp(v',\beta_P(x)).
    \]
\end{proposition}

\begin{proof}
    \leavevmode
    \begin{enumerate}
        \item \textbf{Check typing.}
        If
        \[
            a':u'\to v',
        \]
        then
        \[
            F_1(a'):F_0(u')\to F_0(v').
        \]
        Hence \(\delta_P(F_1a',x)\) is defined whenever
        \(p_P(x)=F_0(v')\), and lies over \(F_0(u')\).

        \item \textbf{Verify the Mealy laws.}
        The unit and multiplication laws follow from functoriality of \(F\) and
        the Mealy laws in \(P\). The output laws are inherited in the same way.

        \item \textbf{Identify residual semantics.}
        Residual execution in \(F^*P\) along a residual arrow \(r'\) is
        execution in \(P\) along \(F_1(r')\). Thus the residual behaviour of
        \((v',x)\) is the residual behaviour of \(x\) precomposed with
        \[
            F_{/v'}:\mathbb A'/v'\to\mathbb A/F_0(v').
        \]
        Hence
        \[
            \beta_{F^*P}(v',x)
            =
            F^\sharp(v',\beta_P(x)).
        \]

        \item \textbf{Display the semantic triangle.}
        This identity is represented by
        \[
        \begin{tikzcd}[column sep=large, row sep=large]
            F^*P
                \arrow[r]
                \arrow[dr, "\beta_{F^*P}"']
            &
            A'_0\times_{A_0}\mathbf{Beh}_{\mathbb A,\mathbb B,0}
                \arrow[d, "F^\sharp"]
            \\
            &
            \mathbf{Beh}_{\mathbb A',\mathbb B,0}.
        \end{tikzcd}
        \]
    \end{enumerate}
\end{proof}

\subsection{Global coequational theories}
\label{subsec:global-coequational-theories}

\begin{definition}[Global coequational theory]
\label{def:global-mealy-coequational-theory}
    A \textbf{global coequational theory} consists of a local coequational theory
    \[
        V_{\mathbb A}\hookrightarrow
        \mathbf{Beh}^{\mathrm{can}}_{\mathbb A,\mathbb B}
    \]
    for each \(\mathbb A\in\mathsf{Inp}\), such that for every internal functor
    \[
        F:\mathbb A'\to\mathbb A
    \]
    the reindexing map
    \[
        F^\sharp:
        A'_0\times_{F_0,A_0,p_{\mathbf{Beh}}}
        \mathbf{Beh}_{\mathbb A,\mathbb B,0}
        \to
        \mathbf{Beh}_{\mathbb A',\mathbb B,0}
    \]
    sends
    \[
        A'_0\times_{F_0,A_0,p_{V_{\mathbb A}}}V_{\mathbb A}
    \]
    into
    \[
        V_{\mathbb A'}.
    \]
    If \(\mathbb B\) is \(\mathbb T\)-structured, a \textbf{global
    \(\mathbb T\)-coequational theory} is such a family for which each
    \(V_{\mathbb A}\) is a local \(\mathbb T\)-coequational theory.
\end{definition}

The global reindexing condition is displayed by
\[
\begin{tikzcd}[column sep=large, row sep=large]
    A'_0\times_{A_0}V_{\mathbb A}
        \arrow[rr, "{F^\sharp|_{V_{\mathbb A}}}"]
        \arrow[dr, dashed]
    &&
    \mathbf{Beh}_{\mathbb A',\mathbb B,0}
    \\
    &
    V_{\mathbb A'}
        \arrow[ur, hook]
    &
\end{tikzcd}
\]
where the fibre product is taken along
\[
    F_0:A'_0\to A_0
    \qquad\text{and}\qquad
    p_{V_{\mathbb A}}:V_{\mathbb A}\to A_0.
\]

\begin{definition}[Global models]
\label{def:global-models}
    Let
    \[
        V=(V_{\mathbb A})_{\mathbb A\in\mathsf{Inp}}
    \]
    be a global coequational theory. Define
    \[
        \mathsf{Mod}^{g}(V)_{\mathbb A}
        =
        \{P\in\mathsf{Rec}_{\mathbb A,\mathbb B}
          \mid
          \beta_P\text{ factors through }V_{\mathbb A}\}.
    \]
\end{definition}

\begin{theorem}[Global coequational closure]
\label{thm:global-coequational-closure}
    Let
    \[
        V=(V_{\mathbb A})_{\mathbb A\in\mathsf{Inp}}
    \]
    be a global coequational theory. Then
    \[
        \mathsf{Mod}^{g}(V)
    \]
    is closed under strict homomorphic preimages, regular quotient recognizers,
    small coproducts, and input pullback along functors in \(\mathsf{Inp}\).

    If \(V\) is a global \(\mathbb T\)-coequational theory, then
    \[
        \mathsf{Mod}^{g}(V)
    \]
    is also closed under \(\mathbb T\)-products.
\end{theorem}

\begin{proof}
    \leavevmode
    \begin{enumerate}
        \item \textbf{Apply local closure pointwise.}
        Closure under strict homomorphic preimages, regular quotient
        recognizers, and small coproducts is
        Theorem~\ref{thm:local-coequational-closure} applied to each
        \(V_{\mathbb A}\). In the structured case, closure under
        \(\mathbb T\)-products is Theorem~\ref{thm:T-coequational-closure}
        applied pointwise.

        \item \textbf{Check input pullback.}
        Let
        \[
            F:\mathbb A'\to\mathbb A
        \]
        and let
        \[
            P\in\mathsf{Mod}^{g}(V)_{\mathbb A}.
        \]
        Choose a factorization
        \[
            \beta_P:P\to V_{\mathbb A}.
        \]
        Proposition~\ref{prop:input-pullback-recognizer} gives
        \[
            \beta_{F^*P}=F^\sharp\beta_P.
        \]
        Since \(V\) is global, \(F^\sharp\) sends the pullback of
        \(V_{\mathbb A}\) into \(V_{\mathbb A'}\). Therefore
        \[
            \beta_{F^*P}
        \]
        factors through \(V_{\mathbb A'}\), so
        \[
            F^*P\in\mathsf{Mod}^{g}(V)_{\mathbb A'}.
        \]
    \end{enumerate}
\end{proof}

\begin{theorem}[Global coequational Birkhoff theorem]
\label{thm:global-coequational-birkhoff}
    Let
    \[
        \mathcal C=(\mathcal C_{\mathbb A})_{\mathbb A\in\mathsf{Inp}}
    \]
    be a replete global family of full recognizer subcategories. Then
    \(\mathcal C\) is of the form
    \[
        \mathcal C_{\mathbb A}
        =
        \mathsf{Mod}^{g}(V)_{\mathbb A}
    \]
    for a global coequational theory
    \[
        V=(V_{\mathbb A})_{\mathbb A\in\mathsf{Inp}}
    \]
    if and only if each \(\mathcal C_{\mathbb A}\) is closed under:
    \begin{enumerate}
        \item strict homomorphic preimages;
        \item regular quotient recognizers;
        \item small coproducts;
    \end{enumerate}
    and the family \(\mathcal C\) is closed under input pullback along functors
    in \(\mathsf{Inp}\).

    If \(\mathbb B\) is \(\mathbb T\)-structured, then global
    \(\mathbb T\)-coequational theories correspond to such families that are
    also closed under \(\mathbb T\)-products.
\end{theorem}

\begin{proof}
    \leavevmode
    \begin{enumerate}
        \item \textbf{Prove necessity.}
        If
        \[
            \mathcal C=\mathsf{Mod}^{g}(V),
        \]
        then the desired closure properties are
        Theorem~\ref{thm:global-coequational-closure}.

        \item \textbf{Construct the local theories from the model family.}
        Conversely, suppose the displayed closure properties hold. For each
        \(\mathbb A\), define
        \[
            V_{\mathcal C,\mathbb A}
            :=
            \bigvee_{P\in\mathcal C_{\mathbb A}}
            \operatorname{Beh}(P)
            \hookrightarrow
            \mathbf{Beh}^{\mathrm{can}}_{\mathbb A,\mathbb B}.
        \]
        By Theorem~\ref{thm:local-coequational-birkhoff},
        \(V_{\mathcal C,\mathbb A}\) is a local coequational theory and
        \[
            \mathcal C_{\mathbb A}
            =
            \mathsf{Mod}(V_{\mathcal C,\mathbb A}).
        \]

        \item \textbf{Check global reindexing on semantic images.}
        Let
        \[
            F:\mathbb A'\to\mathbb A
        \]
        be an internal functor and let
        \[
            P\in\mathcal C_{\mathbb A}.
        \]
        Closure under input pullback gives
        \[
            F^*P\in\mathcal C_{\mathbb A'}.
        \]
        Pulling back the regular epimorphism
        \[
            P\twoheadrightarrow\operatorname{Beh}(P)
        \]
        along
        \[
            A'_0\to A_0
        \]
        gives a regular epimorphism
        \[
            F^*P
            \twoheadrightarrow
            A'_0\times_{A_0}\operatorname{Beh}(P).
        \]
        By Proposition~\ref{prop:input-pullback-recognizer},
        \[
            \beta_{F^*P}=F^\sharp\beta_P.
        \]
        Hence \(F^\sharp\) sends the pullback of \(\operatorname{Beh}(P)\) into
        the semantic image
        \[
            \operatorname{Beh}(F^*P)
            \leq
            V_{\mathcal C,\mathbb A'}.
        \]

        \item \textbf{Descend from the join.}
        The domain
        \[
            A'_0\times_{F_0,A_0,p_{V_{\mathcal C,\mathbb A}}}
            V_{\mathcal C,\mathbb A}
        \]
        is covered by the coproduct of the pullbacks of the semantic images
        \[
            \operatorname{Beh}(P),
            \qquad
            P\in\mathcal C_{\mathbb A}.
        \]
        On each such piece, \(F^\sharp\) lands in
        \(V_{\mathcal C,\mathbb A'}\). By regular-epi descent for subobjects,
        \(F^\sharp\) sends all of \(V_{\mathcal C,\mathbb A}\) into
        \(V_{\mathcal C,\mathbb A'}\). Thus
        \[
            V_{\mathcal C}
            =
            (V_{\mathcal C,\mathbb A})_{\mathbb A\in\mathsf{Inp}}
        \]
        is a global coequational theory.

        \item \textbf{Handle the structured case.}
        If \(\mathbb B\) is \(\mathbb T\)-structured and the family
        \(\mathcal C\) is closed under \(\mathbb T\)-products, then
        Theorem~\ref{thm:T-coequational-birkhoff} applied pointwise shows that
        each \(V_{\mathcal C,\mathbb A}\) is a local
        \(\mathbb T\)-coequational theory. Thus \(V_{\mathcal C}\) is a global
        \(\mathbb T\)-coequational theory.
    \end{enumerate}
\end{proof}

\subsection{Single-object and stateless specializations}
\label{subsec:global-single-object-stateless}

In the one-object case, an internal functor
\[
    F:\mathbb A'\to\mathbb A
\]
is a monoid morphism
\[
    h:M'\to M.
\]
Reindexing residual behaviours is precomposition by \(h\). For a residual
cocycle
\[
    \lambda:M\times M\to N,
\]
the reindexed residual cocycle is
\[
    (h^\sharp\lambda)(r',m')
    =
    \lambda(h(r'),h(m')).
\]
For a language-like behaviour \(L:M\to B_0\), this reduces to
\[
    h^\sharp L=L\circ h.
\]
Thus global coequational closure specializes to closure under inverse morphic
preimages.

In the stateless case, input pullback of a functorial recognizer
\[
    G:\mathbb A\to\mathbb B
\]
along
\[
    F:\mathbb A'\to\mathbb A
\]
is the stateless recognizer induced by the composite
\[
    \mathbb A'
    \xrightarrow{F}
    \mathbb A
    \xrightarrow{G}
    \mathbb B.
\]
Thus global closure under input pullback is ordinary closure of stateless
recognition under precomposition of input functors.

\section{Output-comparison Kleisli endomorphisms}
\label{sec:output-comparison-kleisli-endomorphisms}

This section constructs the internal category that receives the residual action
of a coequational theory.

\subsection{The output-comparison Kleisli monad}
\label{subsec:output-comparison-kleisli-monad}

Define
\[
    T_{\mathbb B}:\mathcal E/B_0\to\mathcal E/B_0
\]
by
\[
    T_{\mathbb B}=(t_B)_!(s_B)^*.
\]
Explicitly, for
\[
    q:X\to B_0,
\]
one has
\[
    T_{\mathbb B}X
    =
    X\times_{q,B_0,s_B}B_1
    \xrightarrow{t_B\pi_{B_1}}
    B_0.
\]
A generalized element of \(T_{\mathbb B}X\) is a pair
\[
    (x,\alpha)
\]
where
\[
    \alpha:q(x)\to b
\]
is an arrow of \(\mathbb B\). The unit is
\[
    x\longmapsto (x,1_{q(x)}),
\]
and multiplication composes comparison arrows in \(\mathbb B\).

The defining pullback is
\[
\begin{tikzcd}[column sep=large, row sep=large]
    T_{\mathbb B}X
        \arrow[r]
        \arrow[d]
        \arrow[dr, phantom, "\lrcorner", very near start]
    &
    B_1
        \arrow[d, "s_B"]
    \\
    X
        \arrow[r, "q"']
    &
    B_0 .
\end{tikzcd}
\]

\begin{proposition}[Kleisli normal form]
\label{prop:kleisli-normal-form}
    A Kleisli map
    \[
        X\rightsquigarrow Y
    \]
    for \(T_{\mathbb B}\), where
    \[
        q_X:X\to B_0,
        \qquad
        q_Y:Y\to B_0,
    \]
    is equivalently a pair
    \[
        F:X\to Y,
        \qquad
        \alpha:X\to B_1
    \]
    such that
    \[
        s_B\alpha=q_YF,
        \qquad
        t_B\alpha=q_X.
    \]
    Thus
    \[
        \alpha_x:q_Y(Fx)\to q_X(x).
    \]

    If
    \[
        (F,\alpha):X\rightsquigarrow Y
    \]
    and
    \[
        (G,\gamma):Y\rightsquigarrow Z
    \]
    are Kleisli maps, their composite is
    \[
        (G,\gamma)\odot(F,\alpha)
        =
        (GF,\alpha\star\gamma),
    \]
    where
    \[
        (\alpha\star\gamma)_x
        =
        \alpha_x\circ \gamma_{Fx}.
    \]
\end{proposition}

\begin{proof}
    \leavevmode
    \begin{enumerate}
        \item \textbf{Unpack a Kleisli map.}
        A Kleisli map \(X\rightsquigarrow Y\) is a map
        \[
            X\to T_{\mathbb B}Y
            =
            Y\times_{q_Y,B_0,s_B}B_1
        \]
        in the slice \(\mathcal E/B_0\). Hence it is equivalently a pair
        \[
            F:X\to Y,
            \qquad
            \alpha:X\to B_1
        \]
        with
        \[
            s_B\alpha=q_YF.
        \]
        The condition that the map lies over \(B_0\) says
        \[
            t_B\alpha=q_X.
        \]

        \item \textbf{Compute Kleisli composition.}
        Applying \((F,\alpha)\) first and \((G,\gamma)\) second gives the chain
        \[
        \begin{tikzcd}[column sep=large, row sep=large]
            q_Z(GFx)
                \arrow[r, "\gamma_{Fx}"]
            &
            q_Y(Fx)
                \arrow[r, "\alpha_x"]
            &
            q_X(x).
        \end{tikzcd}
        \]
        Therefore the composite comparison arrow is
        \[
            \alpha_x\circ\gamma_{Fx}.
        \]
    \end{enumerate}
\end{proof}

Thus Kleisli arrows point in the direction of state transformation, while the
comparison arrow in \(\mathbb B\) points from the output of the transformed
state back to the output of the original state.

\subsection{The internal category of fibrewise Kleisli endomorphisms}
\label{subsec:klend-internal-category}

Let
\[
    p:P\to A_0,
    \qquad
    q:P\to B_0
\]
be any object over
\[
    Z=A_0\times B_0.
\]
For \(v\in A_0\), write
\[
    P_v
\]
for the fibre of \(p\) over \(v\).

\begin{definition}[Kleisli endomorphism category]
\label{def:klend-coeq-theory}
    The \textbf{Kleisli endomorphism category}
    \[
        \operatorname{KlEnd}_{\mathbb B}(P)
    \]
    has object object \(A_0\). An arrow
    \[
        v\to u
    \]
    is a Kleisli map
    \[
        P_v\rightsquigarrow P_u
    \]
    for the output-comparison monad \(T_{\mathbb B}\). Composition is Kleisli
    composition.
\end{definition}

We now spell out the internal representation. Put
\[
    R:=A_0\times A_0
\]
and write
\[
    \sigma,\tau:R\to A_0
\]
for the projections
\[
    \sigma(v,u)=v,
    \qquad
    \tau(v,u)=u.
\]
The pair \((v,u)\) is read as the source and target of an arrow
\[
    v\to u.
\]
Define
\[
    P_{\mathrm{src}}
    :=
    R\times_{\sigma,A_0,p}P,
    \qquad
    P_{\mathrm{tgt}}
    :=
    R\times_{\tau,A_0,p}P.
\]
Over \((v,u)\in R\), these are respectively \(P_v\) and \(P_u\):
\[
\begin{tikzcd}[column sep=large, row sep=large]
    P_{\mathrm{src}}
        \arrow[r]
        \arrow[d]
        \arrow[dr, phantom, "\lrcorner", very near start]
    &
    P
        \arrow[d, "p"]
    \\
    R
        \arrow[r, "\sigma"']
    &
    A_0
\end{tikzcd}
\qquad
\begin{tikzcd}[column sep=large, row sep=large]
    P_{\mathrm{tgt}}
        \arrow[r]
        \arrow[d]
        \arrow[dr, phantom, "\lrcorner", very near start]
    &
    P
        \arrow[d, "p"]
    \\
    R
        \arrow[r, "\tau"']
    &
    A_0 .
\end{tikzcd}
\]

The object \(P_{\mathrm{src}}\) maps to \(R\times B_0\) by
\[
    (v,u,x)\longmapsto ((v,u),q(x)).
\]
The fibrewise output-comparison monad over \(R\), denoted
\[
    T^R_{\mathbb B},
\]
is obtained by applying \(T_{\mathbb B}\) in the \(B_0\)-coordinate and leaving
the \(R\)-coordinate fixed. Thus
\[
    T^R_{\mathbb B}(P_{\mathrm{tgt}})
    :=
    P_{\mathrm{tgt}}
    \times_{q,B_0,s_B}
    B_1
\]
maps to \(R\times B_0\) by
\[
    (v,u,y,\alpha)\longmapsto ((v,u),t_B(\alpha)).
\]
Define
\[
    E_R
    :=
    P_{\mathrm{src}}
    \times_{R\times B_0}
    T^R_{\mathbb B}(P_{\mathrm{tgt}}).
\]
A point of \(E_R\) over
\[
    (v,u,x),
    \qquad
    x\in P_v,
\]
is a pair
\[
    y\in P_u,
    \qquad
    \alpha:q(y)\to q(x).
\]

The pullback defining \(E_R\) is
\[
\begin{tikzcd}[column sep=large, row sep=large]
    E_R
        \arrow[r]
        \arrow[d]
        \arrow[dr, phantom, "\lrcorner", very near start]
    &
    T^R_{\mathbb B}(P_{\mathrm{tgt}})
        \arrow[d]
    \\
    P_{\mathrm{src}}
        \arrow[r]
    &
    R\times B_0 .
\end{tikzcd}
\]

Regard \(E_R\) as an object of \(\mathcal E/P_{\mathrm{src}}\). Because
\(\mathcal E\) is locally cartesian closed, the dependent product along
\[
    \pi_{\mathrm{src}}:P_{\mathrm{src}}\to R
\]
exists. Define
\[
    \operatorname{KlEnd}_{\mathbb B}(P)_1
    :=
    \Pi_{\pi_{\mathrm{src}}}(E_R)
\]
as an object of \(\mathcal E/R\). It represents fibrewise Kleisli maps
\[
    P_v\rightsquigarrow P_u.
\]
Thus a point of
\[
    \operatorname{KlEnd}_{\mathbb B}(P)_1
\]
over \((v,u)\) is a section assigning to every \(x\in P_v\) a pair
\[
    \bigl(\varphi_0(x),\alpha^\varphi_x\bigr)
\]
with
\[
    \varphi_0(x)\in P_u,
    \qquad
    \alpha^\varphi_x:q(\varphi_0x)\to q(x).
\]

The dependent product comes equipped with an evaluation map
\[
    \mathrm{ev}:
    \operatorname{KlEnd}_{\mathbb B}(P)_1
    \times_R
    P_{\mathrm{src}}
    \longrightarrow
    E_R.
\]
In generalized elements,
\[
    \mathrm{ev}(\varphi,x)
    =
    \bigl(\varphi_0(x),\alpha^\varphi_x\bigr).
\]

\begin{proposition}[Internal category of Kleisli endomorphisms]
\label{prop:klend-internal-category}
    The object object \(A_0\) and arrow object
    \[
        \operatorname{KlEnd}_{\mathbb B}(P)_1
    \]
    form an internal category
    \[
        \operatorname{KlEnd}_{\mathbb B}(P).
    \]
    Its identities are induced by the unit of \(T_{\mathbb B}\), and its
    composition is induced by Kleisli composition for \(T_{\mathbb B}\).
\end{proposition}

\begin{proof}
    \leavevmode
    \begin{enumerate}
        \item \textbf{Construct identities.}
        The identity at \(v\) is the Kleisli identity
        \[
            P_v\rightsquigarrow P_v,
            \qquad
            x\longmapsto (x,1_{q(x)}).
        \]
        Internally, this is the transpose of the unit of \(T_{\mathbb B}\) over
        the diagonal
        \[
            A_0\to A_0\times A_0.
        \]

        \item \textbf{Construct composition.}
        For composable arrows
        \[
            \varphi:v\to u,
            \qquad
            \psi:u\to w,
        \]
        write their evaluations as
        \[
            \varphi(x)=\bigl(\varphi_0x,\alpha^\varphi_x\bigr),
            \qquad
            \psi(y)=\bigl(\psi_0y,\alpha^\psi_y\bigr).
        \]
        Their Kleisli composite is
        \[
            (\psi\odot\varphi)(x)
            =
            \bigl(
                \psi_0(\varphi_0x),
                \alpha^\varphi_x\circ\alpha^\psi_{\varphi_0x}
            \bigr).
        \]

        \item \textbf{Display the comparison chain.}
        The comparison arrows form the chain
        \[
        \begin{tikzcd}[column sep=large, row sep=large]
            q(\psi_0(\varphi_0x))
                \arrow[r, "{\alpha^\psi_{\varphi_0x}}"]
            &
            q(\varphi_0x)
                \arrow[r, "{\alpha^\varphi_x}"]
            &
            q(x).
        \end{tikzcd}
        \]

        \item \textbf{Transpose the composite section.}
        The displayed formula gives a section of the defining family \(E_R\)
        over \(P_v\) for the pair \((v,w)\). By the dependent-product
        adjunction, this transposes to the internal composition map.

        \item \textbf{Verify the category laws.}
        Associativity and unitality follow from the monad laws for
        \(T_{\mathbb B}\), equivalently from associativity and unitality of
        composition in \(\mathbb B\).
    \end{enumerate}
\end{proof}

When \(P=V\) is a local coequational theory, we write
\[
    \operatorname{KlEnd}_{\mathbb B}(V)
\]
for the corresponding internal category.

\subsection{Single-object and stateless specializations}
\label{subsec:klend-single-object-stateless}

In the one-object case, \(A_0=1\), so
\[
    \operatorname{KlEnd}_{\mathbb B}(V)
\]
is itself a one-object internal category. Its arrow object is the object of
Kleisli endomaps
\[
    V\rightsquigarrow V.
\]
Thus an arrow is a state transformer
\[
    F:V\to V
\]
together with an output comparison
\[
    \alpha_H:q_V(FH)\to q_V(H)
\]
for each behaviour \(H\in V\). Composition is ordinary composition of state
transformers together with multiplication of comparison arrows in the Kleisli
order
\[
    \alpha_H\circ\gamma_{FH}.
\]

In the stateless situation where \(V\) consists of behaviours represented by
internal functors, a Kleisli endomorphism records how one residual functorial
state is transformed into another, together with an output-comparison arrow at
the root object. This is not a vertical behaviour comparison; it compares only
the anchor outputs. In the chaotic-output case, the comparison component is
determined by endpoints, so a Kleisli endomorphism is essentially a fibrewise
state transformation.

\section{Residual Kleisli actions}
\label{sec:residual-kleisli-actions}

This section packages the Mealy laws as functoriality of a fibrewise Kleisli
action.

\subsection{The residual action}
\label{subsec:residual-action}

\begin{theorem}[Mealy recognizers as fibrewise Kleisli actions]
\label{thm:mealy-recognizers-as-kleisli-actions}
    Let
    \[
        A_0\xleftarrow{p}P\xrightarrow{q}B_0
    \]
    be a carrier span. To give a Mealy recognizer structure on \(P\) is
    equivalent to giving an identity-on-objects internal functor
    \[
        \chi_P:
        \mathbb A^{\mathrm{op}}
        \to
        \operatorname{KlEnd}_{\mathbb B}(P).
    \]
    Under this correspondence, an arrow
    \[
        a:u\to v
    \]
    of \(\mathbb A\), regarded as
    \[
        a^{\mathrm{op}}:v\to u
    \]
    in \(\mathbb A^{\mathrm{op}}\), is sent to the Kleisli map
    \[
        P_v\rightsquigarrow P_u,
        \qquad
        x\longmapsto
        \bigl(\delta(a,x),\omega(a,x)\bigr).
    \]
\end{theorem}

\begin{proof}
    \leavevmode
    \begin{enumerate}
        \item \textbf{Construct the action from a Mealy recognizer.}
        Given \(\delta\) and \(\omega\), the typing equations say precisely
        that, for \(a:u\to v\), the assignment
        \[
            x\longmapsto
            \bigl(\delta(a,x),\omega(a,x)\bigr)
        \]
        is a Kleisli map
        \[
            P_v\rightsquigarrow P_u.
        \]
        By the dependent-product representation of
        \(\operatorname{KlEnd}_{\mathbb B}(P)\), this transposes to an arrow map
        \[
            A_1\to \operatorname{KlEnd}_{\mathbb B}(P)_1
        \]
        over \(A_0\times A_0\), where \(A_1\) is regarded as the arrow object of
        \(\mathbb A^{\mathrm{op}}\) by
        \[
            (t_A,s_A):A_1\to A_0\times A_0.
        \]

        \item \textbf{Identify identities.}
        The Mealy unit laws
        \[
            \delta(1_v,x)=x,
            \qquad
            \omega(1_v,x)=1_{q(x)}
        \]
        say exactly that
        \[
            \chi_P(1_v)
        \]
        is the identity Kleisli endomorphism of \(P_v\).

        \item \textbf{Identify composition.}
        Let
        \[
            a:u\to v,
            \qquad
            b:v\to w.
        \]
        In \(\mathbb A^{\mathrm{op}}\), applying \(b^{\mathrm{op}}\) first and
        \(a^{\mathrm{op}}\) second corresponds to \(b\circ a:u\to w\) in
        \(\mathbb A\). The Mealy multiplication laws
        \[
            \delta(b\circ a,x)=\delta(a,\delta(b,x))
        \]
        and
        \[
            \omega(b\circ a,x)
            =
            \omega(b,x)\circ\omega(a,\delta(b,x))
        \]
        are exactly the Kleisli composition equation
        \[
            \chi_P(b\circ a)=\chi_P(a)\odot\chi_P(b).
        \]

        \item \textbf{Recover the Mealy structure from an action.}
        Conversely, given such an internal functor \(\chi_P\), evaluate
        \(\chi_P(a)\) at \(x\in P_v\). The resulting Kleisli normal form gives
        \[
            \delta(a,x)\in P_u
        \]
        and
        \[
            \omega(a,x):q(\delta(a,x))\to q(x).
        \]
        Identity and composition preservation for \(\chi_P\) are precisely the
        Mealy unit and multiplication laws.

        \item \textbf{Conclude equivalence.}
        The two constructions are inverse by the universal property of the
        dependent product defining
        \(\operatorname{KlEnd}_{\mathbb B}(P)_1\).
    \end{enumerate}
\end{proof}

\begin{definition}[Residual Kleisli action of a theory]
\label{def:residual-kleisli-action}
    Let
    \[
        V\hookrightarrow\mathbf{Beh}^{\mathrm{can}}
    \]
    be a local coequational theory. The \textbf{residual Kleisli action} of
    \(V\) is the internal functor
    \[
        \chi_V:
        \mathbb A^{\mathrm{op}}
        \to
        \operatorname{KlEnd}_{\mathbb B}(V)
    \]
    corresponding, under
    Theorem~\ref{thm:mealy-recognizers-as-kleisli-actions}, to the restricted
    canonical recognizer structure on \(V\).
\end{definition}

Explicitly, for an arrow
\[
    a:u\to v
\]
and a state
\[
    H\in V_v,
\]
one has
\[
    \chi_V(a)(H)
    =
    \bigl(\partial_aH,\omega_V(a,H)\bigr).
\]
The arrow map is obtained by transposition from
\[
    A_1\times_{t_A,A_0,p_V}V
    \longrightarrow
    E_R,
    \qquad
    (a,H)\longmapsto
    \bigl(\partial_aH,\omega_V(a,H)\bigr).
\]
This is displayed by
\[
\begin{tikzcd}[column sep=large, row sep=large]
    A_1\times_R V_{\mathrm{src}}
        \arrow[rr]
        \arrow[dr]
    &&
    E_R
        \arrow[dl]
    \\
    &
    V_{\mathrm{src}}
    &
\end{tikzcd}
\qquad
\leadsto
\qquad
\begin{tikzcd}[column sep=large, row sep=large]
    A_1
        \arrow[r, "\chi_{V,1}"]
        \arrow[dr, "{(t_A,s_A)}"']
    &
    \operatorname{KlEnd}_{\mathbb B}(V)_1
        \arrow[d]
    \\
    &
    A_0\times A_0 .
\end{tikzcd}
\]

\subsection{Functoriality}
\label{subsec:residual-action-functoriality}

\begin{proposition}[Functoriality of the residual action]
\label{prop:residual-action-functorial}
    The assignment
    \[
        \chi_V:
        \mathbb A^{\mathrm{op}}
        \to
        \operatorname{KlEnd}_{\mathbb B}(V)
    \]
    is an internal functor.
\end{proposition}

\begin{proof}
    \leavevmode
    \begin{enumerate}
        \item \textbf{Use the Kleisli-action characterization.}
        By definition, \(\chi_V\) is the fibrewise Kleisli action corresponding
        to the restricted recognizer structure on \(V\).

        \item \textbf{Identify the laws.}
        The identity and composition laws for \(\chi_V\) are exactly the Mealy
        unit and multiplication laws for the restricted canonical recognizer.

        \item \textbf{Display the composite comparison.}
        For
        \[
            a:u\to v,
            \qquad
            b:v\to w,
            \qquad
            H\in V_w,
        \]
        the Kleisli composite has comparison chain
        \[
        \begin{tikzcd}[column sep=large, row sep=large]
            q_V(\partial_a\partial_bH)
                \arrow[r, "{\omega_V(a,\partial_bH)}"]
            &
            q_V(\partial_bH)
                \arrow[r, "{\omega_V(b,H)}"]
            &
            q_V(H).
        \end{tikzcd}
        \]
        The composite is
        \[
            \omega_V(b,H)\circ\omega_V(a,\partial_bH),
        \]
        which is precisely
        \[
            \omega_V(b\circ a,H)
        \]
        by the canonical Mealy output law.
    \end{enumerate}
\end{proof}

\subsection{Single-object and stateless specializations}
\label{subsec:action-single-object-stateless}

In the one-object case, \(\chi_V\) is a one-object internal functor out of
\[
    \mathbb A^{\mathrm{op}}.
\]
Equivalently, it is a homomorphism from the arrow monoid of
\(\mathbb A^{\mathrm{op}}\) to the Kleisli endomorphism monoid of \(V\). Under
the convention
\[
    b\circ a=ba,
\]
passing to \(\mathbb A^{\mathrm{op}}\) reverses the categorical composition. In
the free word-monoid case this makes the arrow monoid of
\[
    \mathbb A_\Sigma^{\mathrm{op}}
\]
carry the usual concatenation order.

The action sends \(m\) to the Kleisli endomorphism
\[
    H\longmapsto
    \bigl(\partial_mH,\omega_V(m,H)\bigr),
\]
where
\[
    (\partial_mH)(r)=H(mr).
\]

In the stateless case induced by an internal functor
\[
    F:\mathbb A\to\mathbb B,
\]
the residual action on the semantic image of \(F\) records how precomposition
by arrows of \(\mathbb A\) transforms the residual functors
\[
    H^F_v:\mathbb A/v\to\mathbb B.
\]
The output comparison associated to
\[
    a:u\to v
\]
is the arrow
\[
    F_1(a):F_0(u)\to F_0(v),
\]
viewed as the canonical comparison between the root outputs of
\[
    H^F_u=\partial_aH^F_v
\]
and \(H^F_v\).

\section{Transition--output images}
\label{sec:transition-output-images}

This section defines the operation image of a local coequational theory and
identifies it as a quotient of \(\mathbb A^{\mathrm{op}}\).

\subsection{Identity-on-objects regular images}
\label{subsec:identity-on-objects-regular-images}

\begin{definition}[Regular epimorphism of identity-on-objects internal categories]
\label{def:regular-epi-identity-on-objects-categories}
    An identity-on-objects internal functor
    \[
        F:\mathbb C\to\mathbb D
    \]
    is a \textbf{regular epimorphism} if its arrow map is a regular epimorphism
    in the slice
    \[
        \mathcal E/(C_0\times C_0),
    \]
    where \(C_0=D_0\). A monomorphism of identity-on-objects internal categories
    is defined similarly by requiring the arrow map to be monic in the same
    slice.
\end{definition}

\begin{lemma}[Identity-on-objects regular images]
\label{lem:identity-on-objects-regular-images}
    Let
    \[
        F:\mathbb C\to\mathbb D
    \]
    be an internal functor whose object map is the identity on a common object
    object \(C_0=D_0\). Factor the arrow map in
    \[
        \mathcal E/(C_0\times C_0)
    \]
    as
    \[
        C_1\xrightarrow{e}I_1\xrightarrow{m}D_1,
    \]
    where \(e\) is a regular epimorphism and \(m\) is a monomorphism. Then
    \(I_1\), with object object \(C_0\), carries a unique internal category
    structure for which
    \[
        \mathbb C\twoheadrightarrow I\hookrightarrow \mathbb D
    \]
    are internal functors.
\end{lemma}

\begin{proof}
    \leavevmode
    \begin{enumerate}
        \item \textbf{Inherit source and target.}
        The source and target maps of \(I_1\) are inherited from its structure
        as an object over \(C_0\times C_0\).

        \item \textbf{Descend identities.}
        Since
        \[
            F_1i_{\mathbb C}=i_{\mathbb D},
        \]
        identities land in the image:
        \[
            i_{\mathbb D}=m(ei_{\mathbb C}).
        \]
        Thus the identity map of \(\mathbb D\) factors through \(I_1\).

        \item \textbf{Construct the regular-epimorphic cover of composable pairs.}
        Let
        \[
            C_2=C_1\times_{C_0}C_1,
            \qquad
            I_2=I_1\times_{C_0}I_1.
        \]
        The induced map
        \[
            C_2\to I_2
        \]
        is a regular epimorphism, since it is built from pullbacks of \(e\) and
        regular epimorphisms are stable under pullback and composition.

        \item \textbf{Descend composition.}
        After precomposition with
        \[
            C_2\twoheadrightarrow I_2,
        \]
        composition in \(\mathbb D\) agrees with the image under \(F\) of
        composition in \(\mathbb C\):
        \[
            m_D(F_1c,F_1d)
            =
            F_1(m_C(c,d))
            =
            m\,e(m_C(c,d)).
        \]
        Thus the composite of two image arrows lands in the image after a
        regular-epimorphic cover. By regular-epi descent for subobjects,
        composition factors uniquely through \(I_1\).

        \item \textbf{Verify the internal-category laws.}
        Unit and associativity laws hold after precomposition with the
        corresponding regular-epimorphic covers from \(\mathbb C\). Hence the
        laws hold in \(I\).
    \end{enumerate}
\end{proof}

\subsection{Operation images}
\label{subsec:operation-images}

\begin{definition}[Transition--output image]
\label{def:operation-image}
    Let
    \[
        V\hookrightarrow\mathbf{Beh}^{\mathrm{can}}
    \]
    be a local coequational theory. The \textbf{transition--output image}, or
    \textbf{operation image}, of \(V\) is the identity-on-objects regular image
    of
    \[
        \chi_V:
        \mathbb A^{\mathrm{op}}
        \to
        \operatorname{KlEnd}_{\mathbb B}(V).
    \]
    It is denoted
    \[
        \operatorname{Op}_{\mathbb A,\mathbb B}(V)
    \]
    or simply
    \[
        \operatorname{Op}(V).
    \]
\end{definition}

Thus \(\operatorname{Op}(V)\) is an internal subcategory of
\[
    \operatorname{KlEnd}_{\mathbb B}(V)
\]
with object object \(A_0\). Its arrow object is the regular image of
\[
    A_1\to \operatorname{KlEnd}_{\mathbb B}(V)_1
\]
in the slice over \(A_0\times A_0\), where \(A_1\), as the arrow object of
\(\mathbb A^{\mathrm{op}}\), lies over \(A_0\times A_0\) by
\[
    (t_A,s_A):A_1\to A_0\times A_0.
\]

The defining image factorization is
\[
\begin{tikzcd}[column sep=large, row sep=large]
    A_1
        \arrow[rr, "\chi_{V,1}"]
        \arrow[dr, two heads]
    &&
    \operatorname{KlEnd}_{\mathbb B}(V)_1
    \\
    &
    \operatorname{Op}(V)_1
        \arrow[ur, hook]
    &
\end{tikzcd}
\qquad
\text{in }\mathcal E/(A_0\times A_0).
\]

\begin{proposition}[Image subcategory]
\label{prop:operation-image-subcategory}
    The transition--output image
    \[
        \operatorname{Op}(V)
    \]
    is an internal subcategory of
    \[
        \operatorname{KlEnd}_{\mathbb B}(V).
    \]
\end{proposition}

\begin{proof}
    \leavevmode
    \begin{enumerate}
        \item \textbf{Apply identity-on-objects image creation.}
        This is Lemma~\ref{lem:identity-on-objects-regular-images} applied to
        the identity-on-objects internal functor
        \[
            \chi_V:
            \mathbb A^{\mathrm{op}}
            \to
            \operatorname{KlEnd}_{\mathbb B}(V).
        \]
    \end{enumerate}
\end{proof}

\subsection{Operation congruences}
\label{subsec:operation-congruences}

\begin{definition}[Operation congruence]
\label{def:operation-congruence}
    For parallel arrows
    \[
        a,a':u\to v
    \]
    of \(\mathbb A\), define
    \[
        a\equiv^{\operatorname{Op}}_V a'
    \]
    if
    \[
        \chi_V(a)=\chi_V(a')
    \]
    as Kleisli maps
    \[
        V_v\rightsquigarrow V_u.
    \]
    Equivalently, for every \(H\in V_v\),
    \[
        \partial_aH=\partial_{a'}H
    \]
    and
    \[
        \omega_V(a,H)=\omega_V(a',H),
    \]
    where the second equality is understood after the first equality identifies
    the sources of the output arrows.
\end{definition}

Internally, this is the kernel pair of
\[
    \chi_{V,1}:A_1\to\operatorname{KlEnd}_{\mathbb B}(V)_1
\]
in
\[
    \mathcal E/(A_0\times A_0),
\]
with \(A_1\) regarded over \(A_0\times A_0\) by
\[
    (t_A,s_A).
\]

\begin{theorem}[Operation image quotient]
\label{thm:operation-image-quotient}
    There is a canonical isomorphism of internal categories
    \[
        \operatorname{Op}(V)
        \cong
        \mathbb A^{\mathrm{op}}/{\equiv^{\operatorname{Op}}_V},
    \]
    where the quotient is the identity-on-objects quotient obtained by
    quotienting the arrow object in
    \[
        \mathcal E/(A_0\times A_0).
    \]
\end{theorem}

\begin{proof}
    \leavevmode
    \begin{enumerate}
        \item \textbf{Identify the congruence.}
        The operation congruence is the kernel pair, in
        \[
            \mathcal E/(A_0\times A_0),
        \]
        of the arrow map of
        \[
            \chi_V:
            \mathbb A^{\mathrm{op}}
            \to
            \operatorname{KlEnd}_{\mathbb B}(V).
        \]
        Since \(\chi_V\) is an internal functor, this kernel pair is compatible
        with identities and composition.

        \item \textbf{Take the effective quotient.}
        Since \(\mathcal E\) is Barr-exact, the slice
        \[
            \mathcal E/(A_0\times A_0)
        \]
        is Barr-exact. Hence the quotient of this kernel pair exists
        effectively and is the regular image of \(\chi_{V,1}\).

        \item \textbf{Identify the quotient with the image.}
        The comparison diagram is
        \[
        \begin{tikzcd}[column sep=large, row sep=large]
            {\equiv^{\operatorname{Op}}_V}
                \arrow[r, shift left=0.5ex]
                \arrow[r, shift right=0.5ex]
            &
            A_1
                \arrow[r, two heads]
                \arrow[dr, two heads]
            &
            A_1/{\equiv^{\operatorname{Op}}_V}
                \arrow[d, dashed, "\cong"]
            \\
            &&
            \operatorname{Op}(V)_1 .
        \end{tikzcd}
        \]
        The internal category structure on the regular image is the one from
        Proposition~\ref{prop:operation-image-subcategory}.
    \end{enumerate}
\end{proof}

\subsection{Contravariance under inclusion}
\label{subsec:operation-image-contravariance}

\begin{proposition}[Contravariance under inclusion]
\label{prop:operation-image-contravariant}
    Let
    \[
        V\hookrightarrow W\hookrightarrow \mathbf{Beh}^{\mathrm{can}}
    \]
    be local coequational theories. Then there is a canonical regular
    epimorphism of internal categories
    \[
        \operatorname{Op}(W)\twoheadrightarrow \operatorname{Op}(V).
    \]
    Thus \(\operatorname{Op}\) is contravariant with respect to inclusion of
    coequational theories: testing equality on the larger theory \(W\) is
    stronger, so the operation congruence for \(W\) is smaller.
\end{proposition}

\begin{proof}
    \leavevmode
    \begin{enumerate}
        \item \textbf{Compare kernels.}
        Suppose
        \[
            a,a':u\to v
        \]
        satisfy
        \[
            \chi_W(a)=\chi_W(a').
        \]
        Then the two operations agree on every state of \(W_v\), hence in
        particular on every state of the subobject
        \[
            V_v\hookrightarrow W_v.
        \]
        Therefore
        \[
            \chi_V(a)=\chi_V(a').
        \]
        Thus
        \[
            \ker(\chi_W)\subseteq \ker(\chi_V)
        \]
        as equivalence relations on \(A_1\) in
        \[
            \mathcal E/(A_0\times A_0).
        \]

        \item \textbf{Descend the quotient.}
        Taking effective identity-on-objects quotients gives a unique internal
        functor
        \[
            \operatorname{Op}(W)\to \operatorname{Op}(V)
        \]
        through which the quotient map
        \[
            A_1\twoheadrightarrow \operatorname{Op}(V)_1
        \]
        factors:
        \[
        \begin{tikzcd}[column sep=large, row sep=large]
            A_1
                \arrow[r, two heads]
                \arrow[dr, two heads]
            &
            \operatorname{Op}(W)_1
                \arrow[d, dashed, two heads]
            \\
            &
            \operatorname{Op}(V)_1 .
        \end{tikzcd}
        \]

        \item \textbf{Show regular epimorphy.}
        The induced arrow map is a regular epimorphism because
        \(\operatorname{Op}(V)_1\) is the effective quotient of
        \(\operatorname{Op}(W)_1\) by the equivalence relation induced from
        \(\ker(\chi_V)\). The object map is the identity on \(A_0\).
    \end{enumerate}
\end{proof}

\subsection{Single-object and stateless specializations}
\label{subsec:operation-single-object-stateless}

In the one-object case, \(\operatorname{Op}(V)\) is a quotient monoid of the
arrow monoid of \(\mathbb A^{\mathrm{op}}\). Two elements
\[
    m,m'\in M
\]
are identified precisely when they induce the same Kleisli endomorphism of
\(V\), equivalently when for every \(H\in V\),
\[
    \partial_mH=\partial_{m'}H
\]
and
\[
    \omega_V(m,H)=\omega_V(m',H)
\]
after identifying the sources of the output arrows.

In the stateless case, \(\operatorname{Op}(V)\) records the image of
\(\mathbb A^{\mathrm{op}}\) acting on the residual functorial behaviours
contained in \(V\). Thus it is the quotient of the input category by the
relation that identifies two arrows when they have the same effect on all
stateless residual behaviours under consideration, including the same root
output comparison.

\section{Syntactic transition--output categories}
\label{sec:syntactic-transition-output-categories}

This section applies the operation-image construction to derivative theories
generated by specifications.

\subsection{Derivative theories generated by specifications}
\label{subsec:derivative-theories-specifications}

\begin{definition}[Derivative theory generated by a specification]
\label{def:derivative-theory-generated-by-specification}
    Let
    \[
        k:K\to\mathbf{Beh}_0
    \]
    be a behaviour specification. Its \textbf{derivative theory} is the
    derivative image
    \[
        \operatorname{Der}_0(K)\hookrightarrow\mathbf{Beh}_0
    \]
    constructed in Chapter~\ref{ch:behaviour-categories-myhill-nerode}. Thus it
    is the regular image of the formal derivative orbit
    \[
        D_K
        =
        A_1\times_{t_A,A_0,p_K}K
        \to
        \mathbf{Beh}_0,
    \]
    where
    \[
        (r,\kappa)\longmapsto \partial_r k(\kappa).
    \]
    The domain \(D_K\) is regarded over \(Z=A_0\times B_0\) by the derivative
    state
    \[
        (r,\kappa)\longmapsto
        \bigl(s_A(r),q_{\mathbf{Beh}}(\partial_r k(\kappa))\bigr).
    \]
\end{definition}

The image factorization is
\[
\begin{tikzcd}[column sep=large, row sep=large]
    D_K
        \arrow[rr, "d_K"]
        \arrow[dr, two heads, "e_K"']
    &&
    \mathbf{Beh}_0
    \\
    &
    \operatorname{Der}_0(K)
        \arrow[ur, hook, "j_K"']
    &
\end{tikzcd}
\qquad
\text{in }\mathcal E/Z.
\]
The derivative-image theorem of
Chapter~\ref{ch:behaviour-categories-myhill-nerode} says that
\[
    \operatorname{Der}_0(K)
\]
is the smallest local coequational theory containing the specified behaviours.

\subsection{Syntactic transition--output categories}
\label{subsec:syntactic-to-categories}

\begin{definition}[Syntactic transition--output category]
\label{def:syntactic-transition-output-category}
    Let
    \[
        k:K\to\mathbf{Beh}_0
    \]
    be a behaviour specification. Its \textbf{syntactic transition--output
    category} is
    \[
        \operatorname{SynTO}_{\mathbb A,\mathbb B}(K)
        :=
        \operatorname{Op}_{\mathbb A,\mathbb B}(\operatorname{Der}_0(K)).
    \]
    Equivalently,
    \[
        \operatorname{SynTO}_{\mathbb A,\mathbb B}(K)
        =
        \operatorname{RegIm}
        \left(
            \mathbb A^{\mathrm{op}}
            \to
            \operatorname{KlEnd}_{\mathbb B}(\operatorname{Der}_0(K))
        \right),
    \]
    where the regular image is taken on arrow objects over
    \[
        A_0\times A_0
    \]
    and the object object remains \(A_0\).
\end{definition}

\begin{definition}[Syntactic operation congruence]
\label{def:syntactic-operation-congruence}
    For parallel arrows
    \[
        a,a':u\to v
    \]
    of \(\mathbb A\), write
    \[
        a\equiv^{\mathrm{syn}}_K a'
    \]
    for the operation congruence
    \[
        a\equiv^{\operatorname{Op}}_{\operatorname{Der}_0(K)}a'.
    \]
    Equivalently, \(a\equiv^{\mathrm{syn}}_K a'\) if \(a\) and \(a'\) induce
    the same Kleisli transition--output operation on the derivative theory
    \[
        \operatorname{Der}_0(K).
    \]
\end{definition}

\subsection{Operation images of recognizers}
\label{subsec:operation-images-recognizers}

\begin{definition}[Operation image of a recognizer]
\label{def:operation-image-recognizer}
    For a recognizer \(P\), define its derivative operation image by
    \[
        \operatorname{Op}(P)
        :=
        \operatorname{Op}(\operatorname{Beh}(P)).
    \]
    Thus \(\operatorname{Op}(P)\) depends only on the semantic image of \(P\).
\end{definition}

\begin{theorem}[Recognizer operation quotient]
\label{thm:recognizer-operation-quotient}
    Let
    \[
        i:K\to P
    \]
    be a realization of the specification
    \[
        k:K\to\mathbf{Beh}_0,
    \]
    so that
    \[
        \beta_P i=k.
    \]
    Then there is a canonical regular epimorphism
    \[
        \operatorname{Op}(P)
        \twoheadrightarrow
        \operatorname{SynTO}_{\mathbb A,\mathbb B}(K).
    \]
    If \(P\) is reachable as a realization of \(K\), then this regular
    epimorphism is an isomorphism.
\end{theorem}

\begin{proof}
    \leavevmode
    \begin{enumerate}
        \item \textbf{Compare semantic images.}
        Since
        \[
            k=\beta_P i,
        \]
        the specified behaviours lie in the semantic image
        \[
            \operatorname{Beh}(P)\hookrightarrow\mathbf{Beh}^{\mathrm{can}}.
        \]
        Since \(\operatorname{Beh}(P)\) is a local coequational theory, the
        smallest local coequational theory containing the specification factors
        through it:
        \[
            \operatorname{Der}_0(K)\hookrightarrow \operatorname{Beh}(P).
        \]

        \item \textbf{Apply contravariance of operation images.}
        By Proposition~\ref{prop:operation-image-contravariant}, this inclusion
        gives a regular epimorphism
        \[
            \operatorname{Op}(\operatorname{Beh}(P))
            \twoheadrightarrow
            \operatorname{Op}(\operatorname{Der}_0(K)).
        \]
        By definition, this is
        \[
            \operatorname{Op}(P)
            \twoheadrightarrow
            \operatorname{SynTO}_{\mathbb A,\mathbb B}(K).
        \]

        \item \textbf{Use reachability for the isomorphism statement.}
        If \(P\) is reachable, then the reachability map
        \[
            \rho_{P,K}:D_K\to P
        \]
        is a regular epimorphism and satisfies
        \[
            \beta_P\rho_{P,K}=d_K.
        \]
        Therefore the regular image of \(\beta_P\) is, by uniqueness of regular
        images in \(\mathcal E/Z\), canonically isomorphic to the regular image
        of \(d_K\):
        \[
        \begin{tikzcd}[column sep=large, row sep=large]
            D_K
                \arrow[r, two heads, "\rho_{P,K}"]
                \arrow[d, two heads, "e_K"']
            &
            P
                \arrow[d, two heads, "e_P"]
            \\
            \operatorname{Der}_0(K)
                \arrow[r, dashed, "\cong"']
            &
            \operatorname{Beh}(P).
        \end{tikzcd}
        \]
        Hence
        \[
            \operatorname{Beh}(P)\cong \operatorname{Der}_0(K)
        \]
        as restricted canonical recognizers. Consequently
        \[
            \operatorname{Op}(P)
            \cong
            \operatorname{SynTO}_{\mathbb A,\mathbb B}(K).
        \]
    \end{enumerate}
\end{proof}

\begin{remark}[Separated realizations]
\label{rem:separated-realizations-operation-image}
    If \(P\) is also behaviourally separated, then the minimal-realization
    theorem of Chapter~\ref{ch:behaviour-categories-myhill-nerode} gives
    \[
        P\cong \mathsf{Min}(K)
    \]
    as a realization. This stronger conclusion is not needed for the operation
    image isomorphism above: reachability alone already identifies the semantic
    image of \(P\) with \(\operatorname{Der}_0(K)\).
\end{remark}

\subsection{Context expansion}
\label{subsec:context-expansion}

\begin{theorem}[Context expansion]
\label{thm:context-expansion}
    Let
    \[
        a,a':u\to v
    \]
    be parallel arrows of \(\mathbb A\). Then
    \[
        a\equiv^{\mathrm{syn}}_K a'
    \]
    if and only if, for every generalized element \(\kappa\in K\) and every
    residual context
    \[
        r:v\to p_K(\kappa),
    \]
    one has
    \[
        \partial_{r\circ a}k(\kappa)
        =
        \partial_{r\circ a'}k(\kappa),
    \]
    and
    \[
        k(\kappa)(a,r)=k(\kappa)(a',r).
    \]
    Here \(k(\kappa)(a,r)\) denotes the arrow assigned by \(k(\kappa)\) to the
    residual morphism
    \[
        r\circ a\to r
    \]
    in \(\mathbb A/p_K(\kappa)\). The second equality is understood after the
    first equality identifies the sources of the output arrows. The first
    equality is an equality in \(\mathbf{Beh}_0\), and hence includes equality
    of the corresponding \(B_0\)-coordinates.
\end{theorem}

\begin{proof}
    \leavevmode
    \begin{enumerate}
        \item \textbf{Cover the generated fibre.}
        The derivative image
        \[
            \operatorname{Der}_0(K)
        \]
        is the regular image of
        \[
            D_K=A_1\times_{t_A,A_0,p_K}K
            \to
            \mathbf{Beh}_0.
        \]
        Pulling back the regular epimorphism
        \[
            D_K\twoheadrightarrow \operatorname{Der}_0(K)
        \]
        along the fibre inclusion
        \[
            \operatorname{Der}_0(K)_v
            \hookrightarrow
            \operatorname{Der}_0(K)
        \]
        gives a regular-epimorphic cover
        \[
            D_K\times_{\operatorname{Der}_0(K)}
            \operatorname{Der}_0(K)_v
            \twoheadrightarrow
            \operatorname{Der}_0(K)_v.
        \]
        Its generalized elements are pairs
        \[
            (r,\kappa),
            \qquad
            r:v\to p_K(\kappa),
        \]
        representing the generated residual state
        \[
            \partial_r k(\kappa).
        \]

        \item \textbf{Display the cover.}
        The cover is displayed by
        \[
        \begin{tikzcd}[column sep=large, row sep=large]
            D_K\times_{\operatorname{Der}_0(K)}
            \operatorname{Der}_0(K)_v
                \arrow[r, two heads]
                \arrow[d]
                \arrow[dr, phantom, "\lrcorner", very near start]
            &
            \operatorname{Der}_0(K)_v
                \arrow[d, hook]
            \\
            D_K
                \arrow[r, two heads]
            &
            \operatorname{Der}_0(K).
        \end{tikzcd}
        \]

        \item \textbf{Compute transitions on the cover.}
        Acting by
        \[
            a:u\to v
        \]
        on the generated state \(\partial_r k(\kappa)\) gives
        \[
            \partial_a(\partial_r k(\kappa))
            =
            \partial_{r\circ a}k(\kappa).
        \]
        Similarly,
        \[
            \partial_{a'}(\partial_r k(\kappa))
            =
            \partial_{r\circ a'}k(\kappa).
        \]

        \item \textbf{Compute outputs on the cover.}
        The output comparison produced by the canonical behaviour recognizer at
        the state
        \[
            \partial_r k(\kappa)
        \]
        along \(a\) is the arrow assigned by \(k(\kappa)\) to the residual
        morphism
        \[
            r\circ a\to r
        \]
        in
        \[
            \mathbb A/p_K(\kappa).
        \]
        This arrow is denoted
        \[
            k(\kappa)(a,r).
        \]
        The analogous output for \(a'\) is
        \[
            k(\kappa)(a',r).
        \]

        \item \textbf{Use regular-epimorphic detection of equality.}
        Equality of the two Kleisli maps
        \[
            \operatorname{Der}_0(K)_v\rightsquigarrow
            \operatorname{Der}_0(K)_u
        \]
        can be checked after precomposition with the displayed
        regular-epimorphic cover. Therefore the two operations induced by
        \(a\) and \(a'\) are equal exactly when the displayed transition and
        output equalities hold for all generated residual states.
    \end{enumerate}
\end{proof}

\subsection{Single-object and stateless specializations}
\label{subsec:syntactic-single-object-stateless}

In the one-object case, a specification is a family of behaviours over the
unique input object. Its derivative theory is generated by all left residual
derivatives
\[
    \partial_m k(\kappa).
\]
The syntactic transition--output category is a one-object category, hence a
monoid, obtained as the quotient of the arrow monoid of
\[
    \mathbb A^{\mathrm{op}}
\]
by equality of transition--output operations on all generated derivatives.

In the stateless case, the specification may consist of behaviours arising
from internal functors
\[
    F_i:\mathbb A\to\mathbb B.
\]
The syntactic transition--output category then identifies two arrows of
\(\mathbb A\) precisely when they induce the same transition and root output
comparison on every residual functorial behaviour generated by the
specification.

\section{Operation congruences and algebraic Eilenberg duality}
\label{sec:operation-congruences-eilenberg}

The coequational Birkhoff theorems above classify recognizer classes by
subrecognizers of the terminal behaviour recognizer. We now extract the
corresponding algebraic side.

For a local coequational theory
\[
    V\hookrightarrow\mathbf{Beh}^{\mathrm{can}}_{\mathbb A,\mathbb B},
\]
the transition--output image
\[
    \operatorname{Op}(V)
\]
is an identity-on-objects quotient of
\[
    \mathbb A^{\mathrm{op}}.
\]
However, the bare quotient need not determine \(V\) itself. The correct
algebraic invariant is the operation congruence
\[
    R_V=\ker(\chi_V),
\]
together with the recognition relation it induces on local coequational
theories.

This leads first to a local Galois connection between coequational theories and
identity-on-objects operation congruences. Its fixed points are the
operation-saturated coequational theories and the operation-closed congruences.
We then globalize this correspondence by transporting congruences through
reindexed coequational theories. The resulting global theorem is a
finiteness-free algebraic Eilenberg duality for the transition--output
semantics of this chapter.

\begin{remark}[Why saturation is necessary]
\label{rem:operation-images-do-not-classify-all-theories}
    The assignment
    \[
        V\longmapsto \operatorname{Op}(V)
    \]
    need not be injective. In the one-object Boolean case, the
    derivative-closed singleton theories generated by the empty language and by
    the full language have the same trivial operation image, although they are
    distinct subobjects of the behaviour object. Thus the algebraic side
    classifies coequational theories after passing to their
    operation-saturated reflection.
\end{remark}

\subsection{Operation congruences}
\label{subsec:eilenberg-operation-congruences}

\begin{definition}[Identity-on-objects operation congruence]
\label{def:identity-on-objects-operation-congruence}
    An \textbf{identity-on-objects operation congruence} on
    \(\mathbb A^{\mathrm{op}}\) is an effective internal equivalence relation
    \[
        R\rightrightarrows A_1
    \]
    in the slice
    \[
        \mathcal E/(A_0\times A_0),
    \]
    where \(A_1\) is regarded as the arrow object of
    \(\mathbb A^{\mathrm{op}}\) by
    \[
        (t_A,s_A):A_1\to A_0\times A_0,
    \]
    such that \(R\) is compatible with identities and composition. Equivalently,
    \(R\) is the kernel pair of an identity-on-objects regular quotient
    \[
        q_R:\mathbb A^{\mathrm{op}}\twoheadrightarrow
        \mathbb A^{\mathrm{op}}/R.
    \]
    We write
    \[
        \mathsf{Cong}_{\mathrm{ioo}}(\mathbb A^{\mathrm{op}})
    \]
    for the poset of such congruences, ordered by inclusion.
\end{definition}

Thus
\[
    R\leq S
\]
means that \(R\) identifies no more arrows than \(S\). Equivalently, the
quotient by \(R\) maps regularly epimorphically onto the quotient by \(S\):
\[
    \mathbb A^{\mathrm{op}}/R
    \twoheadrightarrow
    \mathbb A^{\mathrm{op}}/S.
\]

\begin{definition}[Operation congruence of a local theory]
\label{def:operation-congruence-of-local-theory}
    Let
    \[
        V\hookrightarrow
        \mathbf{Beh}^{\mathrm{can}}_{\mathbb A,\mathbb B}
    \]
    be a local coequational theory. Its \textbf{operation congruence}
    \[
        R_V
    \]
    is the kernel pair of the arrow map
    \[
        \chi_{V,1}:A_1\to
        \operatorname{KlEnd}_{\mathbb B}(V)_1
    \]
    in
    \[
        \mathcal E/(A_0\times A_0).
    \]
\end{definition}

For parallel arrows
\[
    a,a':u\to v
\]
of \(\mathbb A\), the relation
\[
    a\,R_V\,a'
\]
means that \(a\) and \(a'\) induce the same fibrewise Kleisli
transition--output operation on \(V_v\). Equivalently, for every
\[
    H\in V_v,
\]
one has
\[
    \partial_aH=\partial_{a'}H
\]
and
\[
    \omega_V(a,H)=\omega_V(a',H),
\]
where the second equality is typed using the first equality.

\begin{proposition}[Operation image as operation quotient]
\label{prop:operation-image-as-operation-quotient-eilenberg}
    There is a canonical isomorphism of identity-on-objects internal categories
    \[
        \operatorname{Op}(V)
        \cong
        \mathbb A^{\mathrm{op}}/R_V.
    \]
\end{proposition}

\begin{proof}
    By definition, \(\operatorname{Op}(V)\) is the identity-on-objects regular
    image of
    \[
        \chi_V:
        \mathbb A^{\mathrm{op}}
        \to
        \operatorname{KlEnd}_{\mathbb B}(V).
    \]
    Its arrow object is therefore the regular image of
    \[
        \chi_{V,1}:A_1\to
        \operatorname{KlEnd}_{\mathbb B}(V)_1
    \]
    in \(\mathcal E/(A_0\times A_0)\). Since
    \[
        R_V=\ker(\chi_{V,1}),
    \]
    exactness identifies this regular image with the effective quotient
    \[
        A_1/R_V.
    \]
    The internal category structure is the one induced by the
    identity-on-objects image construction. Hence
    \[
        \operatorname{Op}(V)
        \cong
        \mathbb A^{\mathrm{op}}/R_V.
    \]
\end{proof}

\subsection{Recognition by operation quotients}
\label{subsec:recognition-operation-quotients}

\begin{definition}[Recognition by an operation quotient]
\label{def:recognition-by-operation-quotient}
    Let
    \[
        q_R:\mathbb A^{\mathrm{op}}\twoheadrightarrow
        \mathbb A^{\mathrm{op}}/R
    \]
    be an identity-on-objects operation quotient. We say that \(q_R\)
    \textbf{recognizes} a local coequational theory
    \[
        V\hookrightarrow\mathbf{Beh}^{\mathrm{can}}
    \]
    if the residual action
    \[
        \chi_V:\mathbb A^{\mathrm{op}}\to
        \operatorname{KlEnd}_{\mathbb B}(V)
    \]
    factors through \(q_R\):
    \[
    \begin{tikzcd}[column sep=large, row sep=large]
        \mathbb A^{\mathrm{op}}
            \arrow[r, two heads, "q_R"]
            \arrow[dr, "\chi_V"']
        &
        \mathbb A^{\mathrm{op}}/R
            \arrow[d, dashed]
        \\
        &
        \operatorname{KlEnd}_{\mathbb B}(V).
    \end{tikzcd}
    \]
\end{definition}

\begin{lemma}[Recognition criterion]
\label{lem:operation-quotient-recognition-criterion}
    An operation quotient
    \[
        q_R:\mathbb A^{\mathrm{op}}\twoheadrightarrow
        \mathbb A^{\mathrm{op}}/R
    \]
    recognizes \(V\) if and only if
    \[
        R\leq R_V.
    \]
\end{lemma}

\begin{proof}
    The action \(\chi_V\) factors through the quotient \(q_R\) precisely when
    every pair of arrows identified by \(R\) has the same image under
    \(\chi_V\). This is exactly the condition that \(R\) be contained in the
    kernel pair of \(\chi_V\), namely
    \[
        R\leq R_V.
    \]
\end{proof}

\begin{definition}[Theory recognized by a congruence]
\label{def:theory-recognized-by-congruence}
    For
    \[
        R\in\mathsf{Cong}_{\mathrm{ioo}}(\mathbb A^{\mathrm{op}}),
    \]
    define
    \[
        V_R
        :=
        \bigvee
        \left\{
            S\hookrightarrow\mathbf{Beh}^{\mathrm{can}}
            \mid
            S \text{ is a local coequational theory and } R\leq R_S
        \right\}.
    \]
    Thus \(V_R\) is the join of all local coequational theories recognized by
    \(q_R\).
\end{definition}

The collection appearing in the join is a set of subobjects of the fixed object
\[
    \mathbf{Beh}_0
\]
in the fixed slice
\[
    \mathcal E/Z,
\]
relative to the chosen universe.

\begin{proposition}[Largest recognized theory]
\label{prop:largest-recognized-theory}
    The subobject \(V_R\) is a local coequational theory, and
    \[
        R\leq R_{V_R}.
    \]
    Hence \(q_R\) recognizes \(V_R\). Moreover, if \(S\) is any local
    coequational theory recognized by \(q_R\), then
    \[
        S\leq V_R.
    \]
\end{proposition}

\begin{proof}
    \leavevmode
    \begin{enumerate}
        \item \textbf{Form the join.}
        The class of local coequational theories is closed under small joins.
        Therefore
        \[
            V_R
        \]
        is a local coequational theory.

        \item \textbf{Represent the join by a regular image.}
        Write \(V_R\) as the regular image of the coproduct of all local
        theories \(S\) satisfying \(R\leq R_S\):
        \[
        \begin{tikzcd}[column sep=large, row sep=large]
            \coprod_S S
                \arrow[rr]
                \arrow[dr, two heads]
            &&
            \mathbf{Beh}^{\mathrm{can}}
            \\
            &
            V_R
                \arrow[ur, hook]
            &
        \end{tikzcd}
        \]
        in the recognizer category, equivalently on carriers in
        \(\mathcal E/Z\).

        \item \textbf{Descend equality of operations.}
        Let
        \[
            a,a':u\to v
        \]
        be a pair of parallel arrows belonging to \(R\). Pull back the regular
        epimorphism
        \[
            \coprod_S S\twoheadrightarrow V_R
        \]
        along the fibre
        \[
            (V_R)_v\hookrightarrow V_R.
        \]
        The resulting map onto \((V_R)_v\) is again a regular epimorphism. On
        each summand \(S\), the arrows \(a\) and \(a'\) induce equal
        transition--output operations because \(R\leq R_S\). Hence the two
        induced Kleisli operations on \((V_R)_v\) agree after precomposition
        with a regular-epimorphic cover. Therefore they agree on \((V_R)_v\).
        This proves
        \[
            R\leq R_{V_R}.
        \]

        \item \textbf{Prove maximality.}
        If \(S\) is any local coequational theory recognized by \(q_R\), then
        \(R\leq R_S\), so \(S\) appears in the defining join of \(V_R\). Hence
        \[
            S\leq V_R.
        \]
    \end{enumerate}
\end{proof}

\subsection{The local algebraic Eilenberg theorem}
\label{subsec:local-algebraic-eilenberg}

\begin{lemma}[Contravariance of operation congruences]
\label{lem:operation-congruences-contravariant}
    If
    \[
        V\hookrightarrow W\hookrightarrow\mathbf{Beh}^{\mathrm{can}}
    \]
    are local coequational theories, then
    \[
        R_W\leq R_V.
    \]
\end{lemma}

\begin{proof}
    Suppose that two parallel arrows of \(\mathbb A\) induce the same
    transition--output operation on \(W\). Then, by restriction along the
    subrecognizer
    \[
        V\hookrightarrow W,
    \]
    they induce the same transition--output operation on \(V\). Hence every
    pair in \(R_W\) lies in \(R_V\), so
    \[
        R_W\leq R_V.
    \]
\end{proof}

\begin{theorem}[Local operation Galois theorem]
\label{thm:local-operation-galois}
    For every local coequational theory \(V\) and every operation congruence
    \(R\), one has
    \[
        V\leq V_R
        \quad\Longleftrightarrow\quad
        R\leq R_V.
    \]
\end{theorem}

\begin{proof}
    \leavevmode
    \begin{enumerate}
        \item \textbf{Prove the forward implication.}
        Suppose
        \[
            V\leq V_R.
        \]
        By Lemma~\ref{lem:operation-congruences-contravariant},
        \[
            R_{V_R}\leq R_V.
        \]
        By Proposition~\ref{prop:largest-recognized-theory},
        \[
            R\leq R_{V_R}.
        \]
        Hence
        \[
            R\leq R_V.
        \]

        \item \textbf{Prove the reverse implication.}
        Suppose
        \[
            R\leq R_V.
        \]
        Then \(V\) is one of the local coequational theories appearing in the
        join defining \(V_R\). Therefore
        \[
            V\leq V_R.
        \]
    \end{enumerate}
\end{proof}

\begin{definition}[Operation-saturated theory and operation-closed congruence]
\label{def:operation-saturated-and-closed}
    A local coequational theory \(V\) is \textbf{operation-saturated} if
    \[
        V=V_{R_V}.
    \]
    An operation congruence \(R\) is \textbf{operation-closed} if
    \[
        R=R_{V_R}.
    \]
\end{definition}

\begin{lemma}[Closure operators]
\label{lem:operation-closure-operators}
    For every local coequational theory \(V\) and every operation congruence
    \(R\), one has
    \[
        V\leq V_{R_V}
    \]
    and
    \[
        R\leq R_{V_R}.
    \]
    Moreover,
    \[
        R_{V_{R_V}}=R_V
    \]
    and
    \[
        V_{R_{V_R}}=V_R.
    \]
\end{lemma}

\begin{proof}
    The first inequality follows from
    Theorem~\ref{thm:local-operation-galois}, applied to the tautological
    inequality
    \[
        R_V\leq R_V.
    \]
    The second inequality is Proposition~\ref{prop:largest-recognized-theory}.

    For the first displayed equality,
    Proposition~\ref{prop:largest-recognized-theory} gives
    \[
        R_V\leq R_{V_{R_V}},
    \]
    while
    \[
        V\leq V_{R_V}
    \]
    and contravariance give
    \[
        R_{V_{R_V}}\leq R_V.
    \]
    Hence
    \[
        R_{V_{R_V}}=R_V.
    \]

    For the second displayed equality, the local Galois theorem gives
    \[
        V_R\leq V_{R_{V_R}}
    \]
    from the tautological inequality
    \[
        R_{V_R}\leq R_{V_R}.
    \]
    Conversely, since
    \[
        R\leq R_{V_R},
    \]
    antitonicity of \(R\mapsto V_R\) gives
    \[
        V_{R_{V_R}}\leq V_R.
    \]
    Thus
    \[
        V_{R_{V_R}}=V_R.
    \]
\end{proof}

\begin{theorem}[Local algebraic Eilenberg theorem]
\label{thm:local-algebraic-eilenberg}
    The assignments
    \[
        V\longmapsto R_V
    \]
    and
    \[
        R\longmapsto V_R
    \]
    restrict to an order-reversing isomorphism between the poset of
    operation-saturated local coequational theories over
    \((\mathbb A,\mathbb B)\) and the poset of operation-closed
    identity-on-objects operation congruences on
    \(\mathbb A^{\mathrm{op}}\).
\end{theorem}

\begin{proof}
    \leavevmode
    \begin{enumerate}
        \item \textbf{Send saturated theories to closed congruences.}
        Let \(V\) be operation-saturated. Then
        \[
            V=V_{R_V}.
        \]
        Applying \(R_{(-)}\) gives
        \[
            R_V=R_{V_{R_V}},
        \]
        so \(R_V\) is operation-closed.

        \item \textbf{Send closed congruences to saturated theories.}
        Let \(R\) be operation-closed. Then
        \[
            R=R_{V_R}.
        \]
        Applying \(V_{(-)}\) gives
        \[
            V_R=V_{R_{V_R}},
        \]
        so \(V_R\) is operation-saturated.

        \item \textbf{Show that the assignments are inverse on fixed points.}
        If \(V\) is operation-saturated, then
        \[
            V\mapsto R_V\mapsto V_{R_V}=V.
        \]
        If \(R\) is operation-closed, then
        \[
            R\mapsto V_R\mapsto R_{V_R}=R.
        \]

        \item \textbf{Check order reversal.}
        If
        \[
            V\leq W,
        \]
        then Lemma~\ref{lem:operation-congruences-contravariant} gives
        \[
            R_W\leq R_V.
        \]
        Similarly, if
        \[
            R\leq S,
        \]
        then every theory recognized by \(S\) is recognized by \(R\), so
        \[
            V_S\leq V_R.
        \]
    \end{enumerate}
\end{proof}

\subsection{Reindexing local theories and operation congruences}
\label{subsec:reindexing-local-theories-operation-congruences}

We now prepare the global form of the algebraic duality. The point is that
input reindexing transports coequational theories directly, while operation
congruences must be transported through the local Galois connection.

Let
\[
    F:\mathbb A'\to\mathbb A
\]
be an internal functor. Recall that \(F\) induces a residual behaviour
reindexing map
\[
    F^\sharp:
    A'_0\times_{F_0,A_0,p_{\mathbf{Beh}}}
    \mathbf{Beh}_{\mathbb A,\mathbb B,0}
    \to
    \mathbf{Beh}_{\mathbb A',\mathbb B,0}.
\]

\begin{definition}[Reindexed local theory]
\label{def:reindexed-local-theory}
    Let
    \[
        V\hookrightarrow
        \mathbf{Beh}^{\mathrm{can}}_{\mathbb A,\mathbb B}
    \]
    be a local coequational theory. Define
    \[
        F^\bullet V
        \hookrightarrow
        \mathbf{Beh}^{\mathrm{can}}_{\mathbb A',\mathbb B}
    \]
    to be the regular image of the restricted reindexing map
    \[
        A'_0\times_{F_0,A_0,p_V}V
        \longrightarrow
        \mathbf{Beh}_{\mathbb A',\mathbb B,0}
    \]
    induced by \(F^\sharp\). The domain is regarded over
    \[
        Z_{\mathbb A'}=A'_0\times B_0
    \]
    by
    \[
        (v',H)\longmapsto (v',q_V(H)).
    \]
\end{definition}

Thus \(F^\bullet V\) is the local coequational theory generated by all residual
behaviours obtained by reindexing behaviours in \(V\) along \(F\).

\begin{lemma}[Reindexed theories are local]
\label{lem:reindexed-theories-local}
    For every local coequational theory \(V\) over \((\mathbb A,\mathbb B)\),
    the subobject \(F^\bullet V\) is a local coequational theory over
    \((\mathbb A',\mathbb B)\).
\end{lemma}

\begin{proof}
    \leavevmode
    \begin{enumerate}
        \item \textbf{Identify the domain as an input pullback recognizer.}
        The object
        \[
            A'_0\times_{F_0,A_0,p_V}V
        \]
        is the carrier of the input pullback recognizer
        \[
            F^*V:\mathbb A'\rightsquigarrow\mathbb B.
        \]

        \item \textbf{Identify the reindexing map as its semantic map.}
        Since \(V\) is a restricted canonical recognizer, its semantic map is
        its inclusion into
        \[
            \mathbf{Beh}^{\mathrm{can}}_{\mathbb A,\mathbb B}.
        \]
        The input-pullback semantics theorem gives
        \[
            \beta_{F^*V}(v',H)=F^\sharp(v',H).
        \]
        Hence the restricted reindexing map is a strict semantic map from
        \(F^*V\) to the canonical recognizer over \(\mathbb A'\).

        \item \textbf{Take the semantic image.}
        By the regular-image theorem for strict recognizer maps, the regular
        image of this semantic map is a subrecognizer of
        \[
            \mathbf{Beh}^{\mathrm{can}}_{\mathbb A',\mathbb B}.
        \]
        This image is precisely \(F^\bullet V\). Therefore \(F^\bullet V\) is a
        local coequational theory.
    \end{enumerate}
\end{proof}

\begin{lemma}[Basic properties of \(F^\bullet\)]
\label{lem:basic-properties-Fbullet}
    The operation
    \[
        V\longmapsto F^\bullet V
    \]
    is monotone:
    \[
        V\leq W
        \quad\Longrightarrow\quad
        F^\bullet V\leq F^\bullet W.
    \]
    Moreover, for a small family
    \[
        (V_i)_{i\in I}
    \]
    of local coequational theories,
    \[
        F^\bullet\left(\bigvee_i V_i\right)
        =
        \bigvee_i F^\bullet V_i.
    \]
\end{lemma}

\begin{proof}
    Monotonicity follows because the restricted reindexing map for \(V\) factors
    through the restricted reindexing map for \(W\) whenever \(V\leq W\), and
    regular images are monotone on subobjects.

    For joins, write the join
    \[
        \bigvee_i V_i
    \]
    as the regular image of
    \[
        \coprod_i V_i\to\mathbf{Beh}_{\mathbb A,\mathbb B,0}.
    \]
    Pulling back along \(A'_0\to A_0\) and applying \(F^\sharp\), the image of
    the resulting coproduct map is the join of the images of the individual
    summands. Hence
    \[
        F^\bullet\left(\bigvee_i V_i\right)
        =
        \bigvee_i F^\bullet V_i.
    \]
\end{proof}

\begin{definition}[Operational inverse image of a congruence]
\label{def:operational-inverse-image-congruence}
    Let
    \[
        R\in
        \mathsf{Cong}_{\mathrm{ioo}}(\mathbb A^{\mathrm{op}}).
    \]
    Its \textbf{operational inverse image} along
    \[
        F:\mathbb A'\to\mathbb A
    \]
    is the operation congruence
    \[
        F^\flat R
        :=
        R_{F^\bullet(V_R)}
    \]
    on
    \[
        (\mathbb A')^{\mathrm{op}}.
    \]
\end{definition}

Thus \(F^\flat R\) identifies two arrows of \(\mathbb A'\) when they induce the
same transition--output operation on the local theory generated by reindexing
the largest \(\mathbb A\)-theory recognized by \(R\).

\begin{lemma}[Monotonicity of operational inverse image]
\label{lem:Fflat-monotone}
    If
    \[
        R\leq S
    \]
    in
    \[
        \mathsf{Cong}_{\mathrm{ioo}}(\mathbb A^{\mathrm{op}}),
    \]
    then
    \[
        F^\flat R\leq F^\flat S
    \]
    in
    \[
        \mathsf{Cong}_{\mathrm{ioo}}((\mathbb A')^{\mathrm{op}}).
    \]
\end{lemma}

\begin{proof}
    If
    \[
        R\leq S,
    \]
    then every local coequational theory recognized by \(S\) is recognized by
    \(R\). Hence
    \[
        V_S\leq V_R.
    \]
    By Lemma~\ref{lem:basic-properties-Fbullet},
    \[
        F^\bullet V_S\leq F^\bullet V_R.
    \]
    Applying contravariance of operation congruences gives
    \[
        R_{F^\bullet V_R}
        \leq
        R_{F^\bullet V_S}.
    \]
    That is,
    \[
        F^\flat R\leq F^\flat S.
    \]
\end{proof}

\begin{remark}[Comparison with the raw pullback congruence]
\label{rem:raw-pullback-versus-operational-pullback}
    There is also a raw pullback congruence
    \[
        F^{-1}R
    \]
    on \((\mathbb A')^{\mathrm{op}}\), defined by
    \[
        a'\,(F^{-1}R)\,b'
        \quad\Longleftrightarrow\quad
        F_1(a')\,R\,F_1(b')
    \]
    for parallel arrows \(a',b'\) of \(\mathbb A'\).

    The raw pullback congruence is generally too small for global algebraic
    duality. The correct congruence is
    \[
        F^\flat R=R_{F^\bullet(V_R)}.
    \]
    Indeed, since
    \[
        R\leq R_{V_R},
    \]
    any pair of arrows whose \(F\)-images are \(R\)-equivalent acts equally on
    every behaviour in \(V_R\), and therefore acts equally after reindexing
    along \(F\). Consequently there is always an inclusion
    \[
        F^{-1}R\leq F^\flat R.
    \]
    Equality requires additional hypotheses ensuring that behaviour reindexing
    reflects the relevant transition--output equalities. No such reflection
    property is available for arbitrary input functors.
\end{remark}

\subsection{The global algebraic Eilenberg theorem}
\label{subsec:global-algebraic-eilenberg}

We now allow the input category to vary in the chosen category
\[
    \mathsf{Inp}
\]
of internal input categories and internal functors.

The preceding global coequational theories can be described using
\(F^\bullet\): a family
\[
    V=(V_{\mathbb A})_{\mathbb A\in\mathsf{Inp}}
\]
is global precisely when, for every
\[
    F:\mathbb A'\to\mathbb A,
\]
one has
\[
    F^\bullet V_{\mathbb A}\leq V_{\mathbb A'}.
\]

\begin{definition}[Global operation congruence theory]
\label{def:global-operation-congruence-theory-corrected}
    A \textbf{global operation congruence theory} consists of an operation
    congruence
    \[
        R_{\mathbb A}\in
        \mathsf{Cong}_{\mathrm{ioo}}(\mathbb A^{\mathrm{op}})
    \]
    for each input category \(\mathbb A\in\mathsf{Inp}\), such that for every
    internal functor
    \[
        F:\mathbb A'\to\mathbb A
    \]
    one has
    \[
        R_{\mathbb A'}\leq F^\flat R_{\mathbb A}.
    \]
\end{definition}

The inequality
\[
    R_{\mathbb A'}\leq F^\flat R_{\mathbb A}
\]
says that every operation equality imposed at \(\mathbb A'\) is valid on the
local theory generated by reindexing the largest \(\mathbb A\)-theory
recognized by \(R_{\mathbb A}\).

\begin{definition}[Global operation-saturation and operation-closure]
\label{def:global-operation-saturation-closure}
    A global coequational theory
    \[
        V=(V_{\mathbb A})_{\mathbb A}
    \]
    is \textbf{operation-saturated} if each local theory \(V_{\mathbb A}\) is
    operation-saturated:
    \[
        V_{\mathbb A}=V_{R_{V_{\mathbb A}}}.
    \]
    A global operation congruence theory
    \[
        R=(R_{\mathbb A})_{\mathbb A}
    \]
    is \textbf{operation-closed} if each local congruence \(R_{\mathbb A}\) is
    operation-closed:
    \[
        R_{\mathbb A}=R_{V_{R_{\mathbb A}}}.
    \]
\end{definition}

\begin{proposition}[From saturated global theories to closed global congruences]
\label{prop:saturated-global-theories-to-closed-global-congruences}
    If
    \[
        V=(V_{\mathbb A})_{\mathbb A\in\mathsf{Inp}}
    \]
    is an operation-saturated global coequational theory, then
    \[
        R^V_{\mathbb A}:=R_{V_{\mathbb A}}
    \]
    defines an operation-closed global operation congruence theory.
\end{proposition}

\begin{proof}
    \leavevmode
    \begin{enumerate}
        \item \textbf{Check local operation-closure.}
        Since each \(V_{\mathbb A}\) is operation-saturated, the local
        algebraic Eilenberg theorem implies that each
        \[
            R_{V_{\mathbb A}}
        \]
        is operation-closed.

        \item \textbf{Check the global congruence condition.}
        Let
        \[
            F:\mathbb A'\to\mathbb A
        \]
        be an internal functor. Since \(V\) is global,
        \[
            F^\bullet V_{\mathbb A}\leq V_{\mathbb A'}.
        \]
        By contravariance of operation congruences,
        \[
            R_{V_{\mathbb A'}}
            \leq
            R_{F^\bullet V_{\mathbb A}}.
        \]
        Since \(V\) is operation-saturated,
        \[
            V_{\mathbb A}=V_{R_{V_{\mathbb A}}}.
        \]
        Therefore
        \[
            R_{F^\bullet V_{\mathbb A}}
            =
            R_{F^\bullet(V_{R_{V_{\mathbb A}}})}
            =
            F^\flat R_{V_{\mathbb A}}.
        \]
        Hence
        \[
            R^V_{\mathbb A'}\leq F^\flat R^V_{\mathbb A}.
        \]
    \end{enumerate}
\end{proof}

\begin{proposition}[From closed global congruences to saturated global theories]
\label{prop:closed-global-congruences-to-saturated-global-theories}
    If
    \[
        R=(R_{\mathbb A})_{\mathbb A\in\mathsf{Inp}}
    \]
    is an operation-closed global operation congruence theory, then
    \[
        V^R_{\mathbb A}:=V_{R_{\mathbb A}}
    \]
    defines an operation-saturated global coequational theory.
\end{proposition}

\begin{proof}
    \leavevmode
    \begin{enumerate}
        \item \textbf{Check local operation-saturation.}
        Since each \(R_{\mathbb A}\) is operation-closed, the local algebraic
        Eilenberg theorem implies that each
        \[
            V_{R_{\mathbb A}}
        \]
        is operation-saturated.

        \item \textbf{Check globality.}
        Let
        \[
            F:\mathbb A'\to\mathbb A
        \]
        be an internal functor. We must prove
        \[
            F^\bullet V^R_{\mathbb A}\leq V^R_{\mathbb A'}.
        \]
        By definition,
        \[
            V^R_{\mathbb A}=V_{R_{\mathbb A}}.
        \]
        The global congruence condition gives
        \[
            R_{\mathbb A'}\leq F^\flat R_{\mathbb A}.
        \]
        Expanding the definition of \(F^\flat\), this says
        \[
            R_{\mathbb A'}
            \leq
            R_{F^\bullet(V_{R_{\mathbb A}})}.
        \]
        By the local operation Galois theorem, this is equivalent to
        \[
            F^\bullet(V_{R_{\mathbb A}})
            \leq
            V_{R_{\mathbb A'}}.
        \]
        Hence
        \[
            F^\bullet V^R_{\mathbb A}\leq V^R_{\mathbb A'}.
        \]
        Therefore \(V^R\) is a global coequational theory.
    \end{enumerate}
\end{proof}

\begin{theorem}[Global algebraic Eilenberg theorem]
\label{thm:global-algebraic-eilenberg}
    The assignments
    \[
        V=(V_{\mathbb A})_{\mathbb A}
        \longmapsto
        R^V=(R_{V_{\mathbb A}})_{\mathbb A}
    \]
    and
    \[
        R=(R_{\mathbb A})_{\mathbb A}
        \longmapsto
        V^R=(V_{R_{\mathbb A}})_{\mathbb A}
    \]
    restrict to an order-reversing isomorphism between:
    \begin{enumerate}
        \item operation-saturated global coequational theories, and
        \item operation-closed global operation congruence theories.
    \end{enumerate}
\end{theorem}

\begin{proof}
    \leavevmode
    \begin{enumerate}
        \item \textbf{Use the global well-definedness results.}
        Propositions~\ref{prop:saturated-global-theories-to-closed-global-congruences}
        and
        \ref{prop:closed-global-congruences-to-saturated-global-theories}
        show that the two assignments are well-defined at the global level.

        \item \textbf{Apply the local fixed-point theorem pointwise.}
        If \(V\) is operation-saturated, then for every input category
        \(\mathbb A\),
        \[
            V_{\mathbb A}=V_{R_{V_{\mathbb A}}}.
        \]
        Hence
        \[
            V\longmapsto R^V\longmapsto V^{R^V}
        \]
        is the identity on operation-saturated global coequational theories.

        If \(R\) is operation-closed, then for every input category
        \(\mathbb A\),
        \[
            R_{\mathbb A}=R_{V_{R_{\mathbb A}}}.
        \]
        Hence
        \[
            R\longmapsto V^R\longmapsto R^{V^R}
        \]
        is the identity on operation-closed global operation congruence
        theories.

        \item \textbf{Check order reversal.}
        Order reversal is inherited pointwise from the local algebraic
        Eilenberg theorem.
    \end{enumerate}
\end{proof}

\begin{remark}[Why \(F^\flat\), not \(F^{-1}\)]
\label{rem:why-Fflat-not-raw-pullback}
    The raw pullback congruence
    \[
        F^{-1}R
    \]
    only records equality of the \(F\)-images of input arrows. By contrast,
    \[
        F^\flat R
        =
        R_{F^\bullet(V_R)}
    \]
    records equality of the induced transition--output operations on all
    behaviours obtained by reindexing the largest local theory recognized by
    \(R\).

    This distinction is necessary because behaviour reindexing preserves
    residual operations but does not generally reflect equality of behaviours.
    Equality after reindexing only tests residual contexts visible through
    \(F\), whereas a congruence such as \(R_{V_R}\) over \(\mathbb A\) tests all
    residual contexts in \(\mathbb A\). Thus the correct algebraic reindexing is
    obtained by transporting local theories through the local Galois connection,
    not by merely pulling back arrow relations.
\end{remark}

\subsection{The syntactic transition--output universal property}
\label{subsec:syntactic-to-universal-property}

\begin{definition}[Recognition of a specification by an operation quotient]
\label{def:operation-quotient-recognizes-specification}
    Let
    \[
        k:K\to\mathbf{Beh}_0
    \]
    be a behaviour specification. An operation quotient
    \[
        q_R:\mathbb A^{\mathrm{op}}\twoheadrightarrow
        \mathbb A^{\mathrm{op}}/R
    \]
    \textbf{recognizes} \(K\) if it recognizes the derivative theory
    \[
        \operatorname{Der}_0(K)\hookrightarrow\mathbf{Beh}_0.
    \]
\end{definition}

Thus recognition of a specification means recognition of all residual
derivatives generated by that specification.

\begin{theorem}[Syntactic transition--output universal property]
\label{thm:syntactic-to-universal-property}
    Let
    \[
        k:K\to\mathbf{Beh}_0
    \]
    be a behaviour specification, and let
    \[
        q_R:\mathbb A^{\mathrm{op}}\twoheadrightarrow
        \mathbb A^{\mathrm{op}}/R
    \]
    be an identity-on-objects operation quotient. The following are equivalent:
    \begin{enumerate}
        \item \(q_R\) recognizes \(K\);

        \item \(q_R\) recognizes the derivative theory
        \[
            \operatorname{Der}_0(K);
        \]

        \item
        \[
            R\leq R_{\operatorname{Der}_0(K)};
        \]

        \item there is a canonical identity-on-objects regular epimorphism
        \[
            \mathbb A^{\mathrm{op}}/R
            \twoheadrightarrow
            \operatorname{SynTO}_{\mathbb A,\mathbb B}(K).
        \]
    \end{enumerate}
\end{theorem}

\begin{proof}
    \leavevmode
    \begin{enumerate}
        \item \textbf{Unpack recognition of the specification.}
        The equivalence of (1) and (2) is the definition of recognition of a
        specification by an operation quotient.

        \item \textbf{Apply the recognition criterion.}
        By Lemma~\ref{lem:operation-quotient-recognition-criterion},
        recognition of
        \[
            \operatorname{Der}_0(K)
        \]
        is equivalent to
        \[
            R\leq R_{\operatorname{Der}_0(K)}.
        \]
        Hence (2) and (3) are equivalent.

        \item \textbf{Identify the syntactic quotient.}
        By definition,
        \[
            \operatorname{SynTO}_{\mathbb A,\mathbb B}(K)
            =
            \operatorname{Op}(\operatorname{Der}_0(K)).
        \]
        By Proposition~\ref{prop:operation-image-as-operation-quotient-eilenberg},
        \[
            \operatorname{Op}(\operatorname{Der}_0(K))
            \cong
            \mathbb A^{\mathrm{op}}/
            R_{\operatorname{Der}_0(K)}.
        \]

        \item \textbf{Compare quotients.}
        Since quotients are ordered by inclusion of kernel congruences, the
        inequality
        \[
            R\leq R_{\operatorname{Der}_0(K)}
        \]
        is equivalent to the existence of a canonical regular epimorphism
        \[
            \mathbb A^{\mathrm{op}}/R
            \twoheadrightarrow
            \mathbb A^{\mathrm{op}}/
            R_{\operatorname{Der}_0(K)}
            \cong
            \operatorname{SynTO}_{\mathbb A,\mathbb B}(K).
        \]
        Thus (3) and (4) are equivalent.
    \end{enumerate}
\end{proof}

\begin{corollary}[Syntactic transition--output quotient]
\label{cor:syntactic-transition-output-quotient}
    The syntactic transition--output category
    \[
        \operatorname{SynTO}_{\mathbb A,\mathbb B}(K)
    \]
    is the coarsest transition--output quotient of
    \[
        \mathbb A^{\mathrm{op}}
    \]
    recognizing the specification \(K\): every operation quotient recognizing
    \(K\) maps canonically and regularly epimorphically onto
    \[
        \operatorname{SynTO}_{\mathbb A,\mathbb B}(K).
    \]
\end{corollary}

\begin{proof}
    This is the fourth equivalent condition of
    Theorem~\ref{thm:syntactic-to-universal-property}.
\end{proof}

The coequational Birkhoff theorem and the algebraic Eilenberg theorem describe
two complementary aspects of the same residual semantics. The Birkhoff theorem
classifies recognizer classes by subrecognizers of the terminal behaviour
recognizer. The algebraic Eilenberg theorem classifies the operation-saturated
part of this coequational semantics by identity-on-objects operation
congruences on the input categories. For a single specification \(K\), the
associated algebraic object is the syntactic transition--output category
\[
    \operatorname{SynTO}_{\mathbb A,\mathbb B}(K),
\]
which is the universal transition--output quotient recognizing \(K\).

\section{Examples and classical specializations}
\label{sec:coeq-examples-classical-specializations}

This final section collects the main special cases. The first recovers the
classical syntactic monoid of a regular language. The second explains the role
of Boolean closure. The third records compatibility with the classical
Eilenberg correspondence.

\subsection{Boolean output}
\label{subsec:boolean-output-example}

\begin{example}[Boolean output]
\label{ex:boolean-output-theory}
    In \(\mathcal E=\mathbf{Set}\), let \(2=\{0,1\}\) be the usual Boolean
    algebra. More generally, if \(\mathcal E\) contains a chosen Boolean algebra
    object \(2\), the chaotic internal category
    \[
        2_{\mathrm{ch}}
    \]
    on \(2\) is a \(\mathbb T_{\mathrm{Bool}}\)-structured output category. The
    Boolean operations
    \[
        \wedge,\vee,\neg,0,1
    \]
    are internal functors
    \[
        2_{\mathrm{ch}}^n\to 2_{\mathrm{ch}}.
    \]
    Since \(2_{\mathrm{ch}}\) is chaotic, output-comparison arrows carry no
    independent data beyond their endpoints.
\end{example}

Boolean operations enter the theory of language varieties through
\(\mathbb T_{\mathrm{Bool}}\)-coequational theories, namely coequational
subrecognizers closed under the pointwise Boolean algebra structure induced by
\[
    2_{\mathrm{ch}}.
\]
A local Boolean coequational theory is a Boolean algebra of languages closed
under residual derivatives. It is a Boolean algebra of regular languages only
after imposing a finite-state condition or explicitly restricting to regular
languages.

\subsection{The one-object free-monoid case}
\label{subsec:free-monoid-specialization}

Work in
\[
    \mathcal E=\mathbf{Set}.
\]
Let
\[
    \mathbb A_\Sigma
\]
be the one-object internal category associated to the free monoid
\[
    \Sigma^*.
\]
With the one-object convention of
Chapter~\ref{ch:spans-internal-categories-mealy}, categorical composition is
written so that
\[
    b\circ a=ba.
\]
Thus the derivative index is placed on the left:
\[
    (\partial_uL)(w)=L(uw).
\]
Although \(\mathbb A_\Sigma\) uses this categorical convention, the functor
\[
    \chi_L:
    \mathbb A_\Sigma^{\mathrm{op}}
    \to
    \operatorname{KlEnd}_{2_{\mathrm{ch}}}(\operatorname{Der}_0(L))
\]
has the usual word order on its quotient arrow monoid. Indeed, the one-object
category \(\mathbb A_\Sigma\) has categorical composition
\[
    m_{\mathbb A_\Sigma}(u,v)=vu.
\]
Passing to the opposite reverses this, so the arrow monoid of
\[
    \mathbb A_\Sigma^{\mathrm{op}}
\]
has multiplication
\[
    m_{\mathbb A_\Sigma^{\mathrm{op}}}(u,v)=uv,
\]
the usual concatenation product.

Let
\[
    \mathbb B=2_{\mathrm{ch}}
\]
be the chaotic category on the Boolean set \(2=\{0,1\}\). A residual behaviour
is then a language
\[
    L:\Sigma^*\to 2.
\]
For a word \(u\in\Sigma^*\),
\[
    (\partial_uL)(w)=L(uw).
\]

\begin{proposition}[Classical syntactic monoid recovery]
\label{prop:classical-syntactic-monoid-recovery}
    Let
    \[
        L\subseteq\Sigma^*
    \]
    be a regular language, regarded as a specification
    \[
        1\to\mathbf{Beh}_{\mathbb A_\Sigma,2_{\mathrm{ch}},0}.
    \]
    Then the arrow monoid of
    \[
        \operatorname{SynTO}_{\mathbb A_\Sigma,2_{\mathrm{ch}}}(L)
        =
        \operatorname{Op}(\operatorname{Der}_0(L))
    \]
    is canonically isomorphic to the classical syntactic monoid of \(L\).
\end{proposition}

\begin{proof}
    \leavevmode
    \begin{enumerate}
        \item \textbf{Remove the output-comparison data.}
        Since
        \[
            2_{\mathrm{ch}}
        \]
        is chaotic, output-comparison arrows carry no independent information
        beyond their endpoints. Therefore equality of transition--output
        operations is equivalent to equality of the induced transformations on
        residual languages.

        \item \textbf{Use context expansion.}
        By Theorem~\ref{thm:context-expansion}, two words
        \[
            u,v\in\Sigma^*
        \]
        induce the same syntactic transition--output operation if and only if,
        for every residual context \(r\in\Sigma^*\),
        \[
            \partial_{ru}L=\partial_{rv}L.
        \]
        Expanding equality of derivatives gives, for every continuation
        \(s\in\Sigma^*\),
        \[
            rus\in L
            \quad\Longleftrightarrow\quad
            rvs\in L.
        \]

        \item \textbf{Identify the congruence.}
        Thus
        \[
            u\equiv v
            \quad\Longleftrightarrow\quad
            \forall r,s\in\Sigma^*,\;
            rus\in L\Longleftrightarrow rvs\in L,
        \]
        which is the classical syntactic congruence of \(L\).

        \item \textbf{Identify the quotient monoid.}
        The arrow monoid of
        \[
            \operatorname{SynTO}_{\mathbb A_\Sigma,2_{\mathrm{ch}}}(L)
        \]
        is the quotient of the arrow monoid of
        \[
            \mathbb A_\Sigma^{\mathrm{op}}
        \]
        by this congruence. Under the word-order convention described above,
        this quotient multiplication is the usual word multiplication. Hence
        the arrow monoid is canonically the classical syntactic monoid of \(L\).
    \end{enumerate}
\end{proof}

\begin{remark}[Boolean closure and varieties]
\label{rem:boolean-closure-varieties}
    The syntactic monoid of an individual language \(L\) is already recovered
    from the derivative theory
    \[
        \operatorname{Der}_0(L).
    \]
    Passing from individual languages to Boolean algebras of languages is the
    additional language-side closure needed for Eilenberg-type correspondences.
\end{remark}

\subsection{Compatibility with the classical Eilenberg correspondence}
\label{subsec:classical-eilenberg-compatibility}

We finally explain how the preceding constructions specialize to the classical
Eilenberg correspondence.

Work in the one-object Boolean setting:
\[
    \mathbb A_\Sigma
    =
    \text{the one-object category on }\Sigma^*,
    \qquad
    \mathbb B=2_{\mathrm{ch}},
\]
where \(\Sigma\) ranges over finite alphabets. The output category
\(2_{\mathrm{ch}}\) is a Boolean-algebra object in internal categories. Hence
local Boolean coequational theories are Boolean algebras of languages closed
under the residual derivatives supplied by the canonical behaviour recognizer.
In the classical regular-language specialization one restricts attention to
Boolean theories consisting of regular languages, or equivalently to the
finite-state part.

In the standard two-sided classical formulation, one additionally requires
closure under both left and right quotients. The operation-image construction
then detects the usual two-sided syntactic congruence through the context
expansion theorem. Thus the following theorem is not a derivation of
Eilenberg's theorem from local derivative closure alone. Rather, it identifies
the present transition--output image construction with the syntactic monoids
appearing in the classical correspondence, after imposing the usual global
Boolean conditions: Boolean closure, two-sided quotient closure, and inverse
morphic preimage closure.

\begin{definition}[Global Boolean regular language theory]
\label{def:global-boolean-regular-language-theory}
    A \textbf{global Boolean regular language theory} is a family
    \[
        \mathcal V=(\mathcal V_\Sigma)_\Sigma
    \]
    such that:
    \begin{enumerate}
        \item each \(\mathcal V_\Sigma\) is a Boolean algebra of regular
        languages over the finite alphabet \(\Sigma\);
        \item each \(\mathcal V_\Sigma\) is closed under left and right
        quotients;
        \item for every monoid homomorphism
        \[
            h:\Gamma^*\to\Sigma^*,
        \]
        inverse image sends
        \[
            \mathcal V_\Sigma
            \quad\text{into}\quad
            \mathcal V_\Gamma.
        \]
    \end{enumerate}
\end{definition}

Given such a theory \(\mathcal V\), define
\[
    \mathsf P_{\mathcal V}
\]
to be the class of finite monoids dividing finite products of syntactic
monoids
\[
    \operatorname{Syn}(L),
    \qquad
    L\in\mathcal V_\Sigma,
\]
as \(\Sigma\) varies.

Conversely, given a pseudovariety of finite monoids \(\mathsf P\), define
\[
    \mathcal V^{\mathsf P}_\Sigma
    =
    \{L\subseteq\Sigma^*
      \mid
      \operatorname{Syn}(L)\in\mathsf P\}.
\]

\begin{lemma}[Syntactic separation]
\label{lem:syntactic-separation}
    Let \(L\subseteq\Sigma^*\) be regular. Every fibre of the syntactic
    morphism of \(L\) is a Boolean combination of two-sided quotients
    \[
        x^{-1}Ly^{-1}.
    \]
\end{lemma}

\begin{proof}
    \leavevmode
    \begin{enumerate}
        \item \textbf{Use the syntactic congruence.}
        Two words lie in the same syntactic congruence class precisely when they
        cannot be separated by any two-sided context
        \[
            x(-)y
        \]
        with respect to membership in \(L\).

        \item \textbf{Separate distinct classes.}
        Since the syntactic monoid is finite, each congruence class is the
        finite intersection of the two-sided quotient conditions that hold on
        that class and the complements of those that fail on that class.

        \item \textbf{Conclude Boolean definability.}
        Hence each fibre of the syntactic morphism is a Boolean combination of
        two-sided quotients of \(L\).
    \end{enumerate}
\end{proof}

\begin{theorem}[Classical Eilenberg correspondence in the Boolean specialization]
\label{thm:classical-eilenberg-recovery}
    In the global Boolean regular-language specialization, the assignments
    \[
        \mathcal V\longmapsto \mathsf P_{\mathcal V}
    \]
    and
    \[
        \mathsf P\longmapsto \mathcal V^{\mathsf P}
    \]
    recover Eilenberg's correspondence between varieties of regular languages
    and pseudovarieties of finite monoids.
\end{theorem}

\begin{proof}
    \leavevmode
    \begin{enumerate}
        \item \textbf{Identify the categorical monoids with the classical ones.}
        By Proposition~\ref{prop:classical-syntactic-monoid-recovery}, the
        transition--output image associated to a regular language \(L\) has
        arrow monoid equal to the classical syntactic monoid of \(L\).

        \item \textbf{Send a language theory to a pseudovariety.}
        Let
        \[
            \mathcal V=(\mathcal V_\Sigma)_\Sigma
        \]
        be a global Boolean regular language theory. By definition,
        \[
            \mathsf P_{\mathcal V}
        \]
        is closed under finite products and divisors. Hence it is a
        pseudovariety of finite monoids.

        \item \textbf{Send a pseudovariety to a language theory.}
        Let
        \[
            \mathsf P
        \]
        be a pseudovariety of finite monoids. For each finite alphabet
        \(\Sigma\), the class
        \[
            \mathcal V^{\mathsf P}_\Sigma
            =
            \{L\subseteq\Sigma^*
              \mid
              \operatorname{Syn}(L)\in\mathsf P\}
        \]
        is closed under Boolean operations, quotients, and inverse morphic
        preimages, since these operations preserve recognizability by monoids
        in \(\mathsf P\) up to finite products and divisors.

        \item \textbf{Show \(\mathsf P_{\mathcal V^{\mathsf P}}=\mathsf P\).}
        The inclusion
        \[
            \mathsf P_{\mathcal V^{\mathsf P}}\subseteq\mathsf P
        \]
        follows from closure of \(\mathsf P\) under finite products and
        divisors. Conversely, the standard syntactic recognition theorem says
        that every finite monoid \(M\in\mathsf P\) divides a finite product of
        syntactic monoids of languages recognized by \(M\). Those languages
        belong to \(\mathcal V^{\mathsf P}\), so
        \[
            M\in\mathsf P_{\mathcal V^{\mathsf P}}.
        \]

        \item \textbf{Show \(\mathcal V^{\mathsf P_{\mathcal V}}=\mathcal V\).}
        The inclusion
        \[
            \mathcal V\subseteq \mathcal V^{\mathsf P_{\mathcal V}}
        \]
        is immediate. For the reverse inclusion, suppose
        \[
            L\in\mathcal V^{\mathsf P_{\mathcal V}}_\Sigma.
        \]
        Then \(\operatorname{Syn}(L)\) divides a finite product of syntactic
        monoids of languages
        \[
            L_i\in\mathcal V_{\Sigma_i}.
        \]
        By Lemma~\ref{lem:syntactic-separation}, fibres of the syntactic
        morphisms of the \(L_i\) are Boolean combinations of two-sided quotients
        of \(L_i\). Since \(\mathcal V\) is closed under two-sided quotients,
        Boolean operations, and inverse morphic preimages, every language
        recognized by a divisor of this finite product belongs to
        \(\mathcal V_\Sigma\). Hence
        \[
            L\in\mathcal V_\Sigma.
        \]

        \item \textbf{Conclude the correspondence.}
        Thus the two assignments are mutually inverse. Under
        Proposition~\ref{prop:classical-syntactic-monoid-recovery}, this is
        exactly the global Boolean regular-language specialization of the
        transition--output image construction.
    \end{enumerate}
\end{proof}

\subsection{Local and global roles}
\label{subsec:local-global-roles}

\begin{remark}[Local and global roles]
\label{rem:local-global-recovery}
    The local coequational theory records residual closure intrinsic to the
    target-to-source Mealy convention used throughout this thesis. The standard
    two-sided quotient closure of classical language varieties is an additional
    global Boolean specialization condition. The two-sided syntactic congruence
    appears at the transition--output image level, where equality is tested
    against all generated residual states and all continuations.
\end{remark}

\section{Summary}
\label{sec:coequations-syntactic-images-summary}

This chapter has reformulated syntactic transition--output images in terms of
coequational subrecognizers.

\[
\begin{array}{c|l}
\text{construction} & \text{role} \\
\hline
\mathbf{Beh}^{\mathrm{can}}_{\mathbb A,\mathbb B}
&
\text{terminal residual semantic recognizer}
\\
V\hookrightarrow \mathbf{Beh}^{\mathrm{can}}
&
\text{local coequational theory}
\\
\beta_P:P\to\mathbf{Beh}^{\mathrm{can}}
&
\text{unique semantic map to the terminal recognizer}
\\
\operatorname{Beh}(P)
&
\text{regular semantic image of a recognizer}
\\
\mathsf{Mod}(V)
&
\text{recognizers whose semantics land in }V
\\
V_{\mathcal C}
=
\bigvee_{P\in\mathcal C}\operatorname{Beh}(P)
&
\text{coequational theory generated by a model class}
\\
\theta_{\mathbf{Beh}}
&
\text{pointwise structured operation on behaviours}
\\
F^\sharp
&
\text{input reindexing of residual behaviours}
\\
T_{\mathbb B}
&
\text{output-comparison Kleisli monad}
\\
\operatorname{KlEnd}_{\mathbb B}(V)
&
\text{fibrewise Kleisli endomorphism category}
\\
\chi_V
&
\text{residual transition--output action}
\\
\operatorname{Op}(V)
&
\text{transition--output image of a theory}
\\
\operatorname{SynTO}(K)
&
\text{syntactic transition--output category of a specification}
\end{array}
\]

The canonical behaviour recognizer
\[
    \mathbf{Beh}^{\mathrm{can}}_{\mathbb A,\mathbb B}:
    \mathbb A\rightsquigarrow\mathbb B
\]
is terminal in the strict recognizer category. A local coequational theory is a
subrecognizer
\[
    V\hookrightarrow\mathbf{Beh}^{\mathrm{can}}_{\mathbb A,\mathbb B}.
\]
A recognizer \(P\) models \(V\) when its unique semantic map
\[
    \beta_P:P\to\mathbf{Beh}^{\mathrm{can}}
\]
factors through \(V\). Semantic images
\[
    \operatorname{Beh}(P)\hookrightarrow\mathbf{Beh}^{\mathrm{can}}
\]
are regular images of these terminal semantic maps.

The local coequational Birkhoff theorem says that, in the Grothendieck topos
setting of this chapter, local coequational theories are equivalent to replete
full recognizer classes closed under strict homomorphic preimages, regular
quotient recognizers, and small coproducts. The theory generated by such a
class \(\mathcal C\) is the join
\[
    V_{\mathcal C}
    =
    \bigvee_{P\in\mathcal C}\operatorname{Beh}(P).
\]
If the output category \(\mathbb B\) is a \(\mathbb T\)-algebra object in
internal categories, then residual behaviours inherit pointwise
\(\mathbb T\)-structure. Local \(\mathbb T\)-coequational theories correspond
to the same model classes with the additional closure condition under
\(\mathbb T\)-products.

Allowing the input category to vary gives global coequational theories, stable
under reindexing of behaviours along input functors. The global coequational
Birkhoff theorem says that these correspond to global model families closed
pointwise under strict homomorphic preimages, regular quotients, and small
coproducts, and globally under input pullback. In the structured case, one adds
closure under \(\mathbb T\)-products.

The output-comparison monad
\[
    T_{\mathbb B}=(t_B)_!(s_B)^*
\]
gives a Kleisli category of output-comparison maps. A Mealy recognizer is
equivalently a fibrewise Kleisli action
\[
    \chi_P:
    \mathbb A^{\mathrm{op}}
    \to
    \operatorname{KlEnd}_{\mathbb B}(P).
\]
For a local coequational theory \(V\), this gives the residual action
\[
    \chi_V:
    \mathbb A^{\mathrm{op}}
    \to
    \operatorname{KlEnd}_{\mathbb B}(V).
\]
Its transition--output image is the identity-on-objects regular image
\[
    \operatorname{Op}(V)
    =
    \operatorname{RegIm}(\chi_V).
\]

For a specification \(K\), the syntactic transition--output category is
\[
    \operatorname{SynTO}_{\mathbb A,\mathbb B}(K)
    =
    \operatorname{Op}(\operatorname{Der}_0(K)).
\]
In the one-object Boolean case, the arrow monoid of
\[
    \operatorname{SynTO}(L)
\]
is the classical syntactic monoid of \(L\). In the global Boolean
regular-language specialization, the construction is compatible with
Eilenberg's correspondence between varieties of regular languages and
pseudovarieties of finite monoids.
